\documentclass[11pt]{article}
\usepackage[margin=1in]{geometry}
\IfFileExists{lmodern.sty}{\usepackage[T1]{fontenc}\usepackage{lmodern}}{\DeclareTextCommand{\k}{OT1}[1]{\c{##1}}}% ogonek fallback when Latin Modern is not installed
\usepackage{amsmath,amssymb,amsthm,mathtools}
\usepackage{stmaryrd}
\usepackage{booktabs}
\usepackage{longtable,array}
\usepackage[hidelinks]{hyperref}
\usepackage[capitalise,nameinlink]{cleveref}
\usepackage{tikz-cd}
\setlength{\emergencystretch}{2em}

% ---------- theorem environments (as in KBFold) ----------
\newtheorem{theorem}{Theorem}[section]
\newtheorem{lemma}[theorem]{Lemma}
\newtheorem{proposition}[theorem]{Proposition}
\newtheorem{corollary}[theorem]{Corollary}
\newtheorem{fact}[theorem]{Fact}
\theoremstyle{definition}
\newtheorem{definition}[theorem]{Definition}
\newtheorem{remark}[theorem]{Remark}
\newtheorem{example}[theorem]{Example}

% ---------- notation (KBFold notation, plus the monomial-basis objects of this paper) ----------
\newcommand{\Fq}{\mathbb{F}_q}
\newcommand{\K}{\mathbb{K}}
\newcommand{\N}{\mathbb{N}}
\newcommand{\iv}[2]{\llbracket #1, #2 \rrbracket}
\newcommand{\wt}{\mathrm{wt}}
\newcommand{\eq}{\mathrm{eq}}
\newcommand{\RS}{\mathrm{RS}}
\newcommand{\Enc}{\mathrm{Enc}}
\newcommand{\fold}{\mathrm{fold}}
\newcommand{\cfold}{\mathrm{cfold}}
\newcommand{\coef}[2]{[X^{#1}]\,#2}
\newcommand{\Cc}{\mathcal{C}}
\newcommand{\Dd}{\mathcal{D}}
\newcommand{\Pp}{\mathcal{P}}
\newcommand{\lo}{\mathrm{lo}}
\newcommand{\res}{\mathrm{res}}
\newcommand{\ev}{\mathrm{ev}}
\newcommand{\F}{\mathbb{F}}
\newcommand{\mlin}{\mathrm{mlin}}
\newcommand{\Pieval}{\Pi_{\mathrm{eval}}}
\newcommand{\Pisum}{\Pi_{\mathrm{sum}}}
\newcommand{\KB}[1]{KBFold~#1}% reference to a numbered result of the companion paper [Mkh26]
\newcommand{\sfversion}{0.1.0}

\title{Sum-FRI: Hypercube Sums and Coefficient-Form Multilinear Evaluations\\ from Shared-Challenge FRI Folding}
\author{Abdelali Mkhida\\ Algorizk Labs\\ \texttt{ali.mkhida@algorizk.xyz}\\ ORCID \href{https://orcid.org/0009-0009-2101-9070}{\texttt{0009-0009-2101-9070}}}
\date{October 2026}

\begin{document}
\maketitle

\begin{abstract}
We prove hypercube sums $\sum_a f(a)$ and coefficient-form multilinear evaluations $P_f(z) = \sum_a f(a) \prod_k z_k^{a_k}$ for a table $f$ committed as the Reed--Solomon codeword of $V_f = \sum_a f(a) X^a$ on a smooth domain.
Classical FRI folding of $V_f$ restricts $P_f$ in its first variable, and $n$ folds give $P_f(\theta_1, \dots, \theta_n)$.
When the folding challenges also drive the claim, each round costs one affine round polynomial, sent as a single field element, with no equality table and no change of basis; for the sum, the round messages are plain sums of the table.
This round structure is that of the evaluation protocol of DeepFold.
We give a complete proof of soundness, evaluation binding and round-by-round knowledge soundness in the unique-decoding regime, whose only external input about codes is the correlated-agreement theorem of Ben-Sasson, Carmon, Ishai, Kopparty and Saraf, and derive post-quantum knowledge soundness of the compiled scheme in the quantum random oracle model from the theorem of Chiesa, Di, Hu and Zheng.
A restriction round costs one bad challenge, against two for a degree-$2$ sumcheck.
The proofs are formalised in Lean~4 on top of the KBFold library, with correlated agreement as the only axiom.
A Rust implementation that shares its field, NTT, Merkle trees and transcript with KBFold opens a sum $27$--$41\%$ faster than KBFold and than a coefficient-form sumcheck scheme, at equal commitment cost, verification time and proof size, and $34$--$57\%$ faster with post-quantum parameters.
\end{abstract}

\tableofcontents

\section{Introduction}\label{sec:intro}

Let $f$ be a table of $N = 2^n$ field elements indexed by the Boolean hypercube $B_n = \{0,1\}^n$.
Two claims about $f$ recur in succinct arguments: the \emph{hypercube sum}
\[
  S \;=\; \sum_{b \in B_n} f(b),
\]
which is the statement of the sumcheck protocol~\cite{LFKN92}, and the \emph{evaluation} of a multilinear polynomial attached to $f$ at a point $z \in \F^n$, which is the statement of a multilinear polynomial commitment scheme.
Hash-based schemes in the BaseFold paradigm~\cite{ZCF24} commit to a Reed--Solomon encoding of $f$ and prove such claims by interleaving a sumcheck with FRI folding~\cite{BBHR18}, the two sharing their random challenges.

This paper studies the simplest member of this family: the table is committed \emph{as the coefficient vector} of a univariate polynomial, and the scalar claim is reduced by a restriction of one variable per round, with no sumcheck.
Put
\[
  V_f(X) \;=\; \sum_{b=0}^{N-1} f(b)\, X^{b},
  \qquad
  P_f(x_1,\dots,x_n) \;=\; \sum_{b \in B_n} f(b) \prod_{k=1}^{n} x_k^{b_k},
\]
so that $P_f$ is the multilinear polynomial whose \emph{monomial coefficients} are the entries of $f$, and $V_f$ is its Kronecker substitution, $V_f(X) = P_f(X, X^2, \dots, X^{2^{n-1}})$.
Then
\[
  \sum_{b \in B_n} f(b) \;=\; P_f(1,\dots,1) \;=\; V_f(1).
\]
The classical FRI fold $\cfold_\theta(U) = U_e + \theta\, U_o$, where $U(X) = U_e(X^2) + X U_o(X^2)$, substitutes $\theta$ for the first variable of the coefficient-form multilinear polynomial (\cref{sec:restriction}).
For a claim $P_f(z) = v$, the round message of the prover is the affine polynomial
\[
  g_j(T) \;=\; P_{f}(\theta_1,\dots,\theta_{j-1},\, T,\, z_{j+1},\dots,z_n),
\]
the verifier checks $g_j(z_j)$ against the current claim, draws $\theta_j$, and the same $\theta_j$ folds the committed codeword.
For the hypercube sum, $z = (1,\dots,1)$ and $g_j(T) = U_e(1) + T\, U_o(1)$, where $U$ is the current folded polynomial.
After $n$ rounds the folded codeword is the constant $P_f(\theta_1,\dots,\theta_n)$, and the verifier compares it with the reduced claim.
Three states thus move under one challenge sequence: the scalar claim, the coefficient-form polynomial, and the Reed--Solomon word.

The round structure is that of the evaluation protocol of DeepFold~\cite{GLHQTZ25}, which proves coefficient-form evaluations by the same restrictions; DeepFold adds out-of-domain points to reach the list-decoding regime.
The present paper gives the unique-decoding version a complete, self-contained treatment, with the hypercube sum as its distinguished instance, and develops it as the coefficient-form companion of the Boolean-kernel scheme KBFold~\cite{Mkh26}, on the same code and with the same proof machinery.

\subsection{Contributions}\label{sec:intro:contrib}

\paragraph{The restriction dictionary (\cref{sec:restriction}).}
We state the monomial-basis counterpart of the KBFold fold dictionary: the classical fold of $V_f$ is the Kronecker substitution of the restriction of $P_f$ in its first variable, and $n$ folds with $\theta_1,\dots,\theta_n$ give the constant $P_f(\theta_1,\dots,\theta_n)$.
For a word $w$ on a smooth domain $L$, the classical word fold is the line $w_e + \theta\, w_o$ whose generators determine $w$.
This is the property through which correlated agreement enters the soundness proof, and it holds for the classical fold without the reparametrisation of the kernel fold (\KB{Proposition~5.11}).
We record the relation between the two encodings of the same sum: $V_f(1)$ in the monomial basis, and $[X^{N-1}]\,U_f$ in the Boolean-kernel basis.

\paragraph{The scheme and its completeness (\cref{sec:scheme}).}
We define the commitment, the evaluation protocol $\Pieval$ for claims $P_f(z) = v$, and its specialisation $\Pisum$ to hypercube sums, with a stopping parameter, committed levels and the local query checks of KBFold.
The prover sends one affine polynomial per round, that is, one field element once the restriction check $s_j(z_j) = v_{j-1}$ is used to recover the other.
It never forms an equality table, and the commitment requires no change of basis: the table is the coefficient vector.

\paragraph{Soundness (\cref{sec:sound}).}
We prove soundness and evaluation binding in the unique-decoding regime, round-by-round knowledge soundness with a decoding extractor, and the relaxed round-by-round knowledge soundness of Chiesa, Di, Hu and Zheng~\cite{CDHZ25}, which yields post-quantum knowledge soundness of the compiled scheme in the quantum random oracle model.
The proof follows KBFold: the codeword chain and the passing sets carry over unchanged, the folding lemmas are re-proved for the classical fold, and the sumcheck step is replaced by a one-line restriction step.
That step costs $1/|\F|$ per round, against $2/|\F|$ for the degree-$2$ sumcheck of KBFold, and in every round it is dominated by the folding error $|L_j|/|\F|$.
The only external result about codes is again the correlated-agreement theorem of Ben-Sasson, Carmon, Ishai, Kopparty and Saraf~\cite{BCIKS23}.

\paragraph{Formal verification (\cref{app:lean}).}
A Lean~4 formalisation, built on the KBFold Lean library and checked against the same single axiom (correlated agreement), covers the restriction dictionary, completeness, soundness, binding, round-by-round knowledge soundness, the combinatorial part of relaxed round-by-round knowledge soundness, and the transfer to committed levels.

\paragraph{Implementation and experiments (\cref{sec:impl,sec:exp}).}
A Rust implementation reuses the field, NTT, salted Merkle trees and transcript of the KBFold crate.
We compare it on the same machine with four alternatives for the same statements:
\begin{itemize}
  \item the KBFold scheme, which proves the sum as $2^n \tilde f(\tfrac12,\dots,\tfrac12)$;
  \item the coefficient-form scheme with a degree-$2$ sumcheck, as in BaseFold and in the baseline of KBFold;
  \item the quotient opening of $V_f$ at $X = 1$ on a coset domain;
  \item WHIR~\cite{ACFY24}.
\end{itemize}
With the same commitment, verification time and proof size, the opening of a sum is $27$--$41\%$ faster than with KBFold or with the coefficient-form sumcheck scheme, and $34$--$57\%$ faster with post-quantum parameters.

\subsection{Positioning}\label{sec:intro:position}

\Cref{tab:position} lists the protocols that prove hypercube sums or coefficient-form evaluations with univariate encodings, and how each one enters this paper.

\begin{table}[t]
\centering
\small
\begin{tabular}{@{}p{0.2\textwidth}p{0.36\textwidth}p{0.36\textwidth}@{}}
\toprule
Work & Mechanism & Role in this paper \\
\midrule
Sumcheck~\cite{LFKN92} & one variable per round, round polynomial of degree $d$, final evaluation query & the reduction that the restriction step replaces for linear statements (\cref{sec:restriction}) \\
Aurora~\cite{BCRSVW19} & univariate sumcheck over a multiplicative subgroup, for evaluation-form polynomials & different encoding; not compared \\
Gemini~\cite{BCHO22}, PolyFRIM~\cite{ZLGSCLD24} & coefficient-form evaluation by folding at the evaluation point $z$; one commitment per fold & same identity $P_f(z)$ by even--odd splitting; our folds use fresh random challenges and need no extra commitment \\
Zeromorph~\cite{KT24} & the table as univariate coefficients, $\sum_b f(b) X^{b}$; quotients and degree checks with KZG & the same encoding $V_f$ with a pairing-based proof; not hash-based \\
BaseFold~\cite{ZCF24} & FRI folds whose last constant is a random evaluation; degree-$2$ sumcheck on $f\cdot\eq_z$ & source of the final-value identity; the coefficient-form sumcheck baseline of \cref{sec:exp} \\
DeepFold~\cite{GLHQTZ25} & coefficient-form evaluation by restrictions sharing the FRI challenges; out-of-domain points for list decoding & the round structure of $\Pieval$; we treat its unique-decoding version, with full proofs \\
FriEval~\cite{Zha25} & FRI as a protocol evaluating a committed multilinear at challenge points & the same reading of FRI, in a functional framework \\
Fold-DCS~\cite{LMN25} & sumcheck in $O(\log n)$ rounds by divide and conquer & different technique; not compared \\
WHIR~\cite{ACFY24}, Hab\"ock~\cite{Hab24} & constrained Reed--Solomon codes; list decoding via subcode correlated agreement & external baseline; route to the list-decoding version (\cref{sec:concl}) \\
KBFold~\cite{Mkh26} & table encoded in the Boolean-kernel basis; kernel fold evaluates $\tilde f$ & companion scheme: same code, soundness machinery and Lean library \\
\bottomrule
\end{tabular}
\caption{Protocols for hypercube sums and coefficient-form evaluations with univariate encodings.}
\label{tab:position}
\end{table}

The coefficient-form polynomial $P_f$ is not the multilinear extension $\tilde f$ of the table: $P_f(b) = \sum_{c \preceq b} f(c)$ on $B_n$.
A system whose statements concern $\tilde f$, such as a sumcheck-based argument that ends with a claim $\tilde f(r) = v$, can use $\Pieval$ after the M\"obius transform of the table, which is what BaseFold and DeepFold do, or use KBFold, which avoids the transform.
The hypercube sum is the one statement on which the two polynomials agree without any transform, since $\sum_b f(b) = P_f(1,\dots,1)$ and, in odd characteristic, $\sum_b f(b) = 2^n \tilde f(\tfrac12,\dots,\tfrac12)$.

\subsection{Related work}\label{sec:intro:related}

\paragraph{Proximity tests and their soundness.}
FRI~\cite{BBHR18} tests proximity to Reed--Solomon codes by repeated even--odd folding, and its soundness rests on proximity gaps for Reed--Solomon codes~\cite{BCIKS23}.
DEEP-FRI~\cite{BGKS20} samples outside the evaluation domain to bind a word to one codeword in the list-decoding regime, the technique DeepFold uses for multilinear evaluations.
WHIR~\cite{ACFY24} tests proximity to constrained Reed--Solomon codes with very fast verification.
Hab\"ock~\cite{Hab24} proves BaseFold sound up to the Johnson bound through correlated agreement for linear subcodes.
Zhang~\cite{Zha25} recasts FRI as a functional proof system evaluating committed multilinear polynomials and builds the security of BaseFold on the theorems of FRI.
Our analysis stays in the unique-decoding regime and uses only the correlated-agreement theorem.

\paragraph{Multilinear commitments from folding.}
BaseFold~\cite{ZCF24} introduced the interleaving of sumcheck rounds with folding rounds under shared challenges and the observation that the last FRI constant is a random evaluation of the committed coefficient-form polynomial.
DeepFold~\cite{GLHQTZ25} replaces the sumcheck by restrictions for coefficient-form polynomials and reaches the list-decoding regime; it also batches polynomials of different sizes and adds zero knowledge.
KBFold~\cite{Mkh26} encodes the table in the Boolean-kernel basis so that the folds evaluate the multilinear extension of the table itself.

\paragraph{Multilinear-to-univariate reductions.}
Gemini~\cite{BCHO22} reduces a coefficient-form claim $P(z) = v$ to univariate claims by folding at the coordinates of $z$ and committing to every folded polynomial; PolyFRIM~\cite{ZLGSCLD24} instantiates this approach with FRI.
Zeromorph~\cite{KT24} maps the table to univariate coefficients, which is the encoding $V_f$ of this paper, and proves evaluations of $\tilde f$ with quotients and degree checks.
Mohamed~\cite{Moh25} classifies these adaptors by how they map multilinear to univariate polynomials and gives a values-to-values reduction.
These works treat the univariate commitment as a black box; here the univariate polynomial is committed through a Reed--Solomon code whose own folding carries the reduction.

\paragraph{Sumcheck variants.}
The sumcheck protocol~\cite{LFKN92} reduces a hypercube sum to one evaluation, with one round polynomial of degree $d$ per variable.
Aurora~\cite{BCRSVW19} proves sums over multiplicative subgroups for univariate polynomials in one round.
Fold-DCS~\cite{LMN25} reduces the number of rounds to $O(\log n)$ by a divide-and-conquer recursion.
For the linear statement $\sum_b f(b) = S$ about a committed coefficient vector, the restriction step of this paper plays the role of the sumcheck, with affine round polynomials and no final evaluation query.

\subsection{Organisation}\label{sec:intro:org}

\Cref{sec:prelim,sec:proofs} recall, without proofs, the notation and the results of KBFold that we use: multilinear polynomials, smooth domains, Reed--Solomon codes, correlated agreement, interactive oracle proofs, round-by-round soundness and the theorem of~\cite{CDHZ25}.
\Cref{sec:restriction} proves the restriction dictionary and the line form of the classical fold.
\Cref{sec:scheme} defines the scheme and proves completeness, and \cref{sec:sound} proves soundness, evaluation binding, round-by-round knowledge soundness and post-quantum knowledge soundness.
\Cref{sec:impl,sec:exp} describe the implementation and the measurements, and \cref{sec:concl} concludes.
\Cref{app:lean} maps the paper to the Lean formalisation.

Results of KBFold are cited by their numbers in~\cite{Mkh26}, as in \KB{Theorem~2.23}.
Every result specific to this paper is proved here.

\section{Preliminaries}\label{sec:prelim}

This section fixes notation and recalls the algebraic facts that we use.
All of them are proved in KBFold~\cite{Mkh26}; we restate them without proof and give the number of each result there.
The only external result about codes is the correlated-agreement theorem (\cref{thm:bciks}).

\subsection{Integers, bits and points}\label{sec:prelim:bits}

$\N = \{0,1,2,\dots\}$, and for integers $a, c$ we write $\iv{a}{c} = \{i \in \mathbb{Z} : a \le i \le c\}$.
The letter $n \in \N$ denotes a number of variables and $N = 2^n$.
Every $b \in \iv{0}{N-1}$ has a unique binary decomposition $b = \sum_{k=1}^n b_k 2^{k-1}$ with $b_k \in \{0,1\}$; bit $b_1$ is the least significant one, and we identify $b$ with its bit vector $(b_1,\dots,b_n) \in B_n = \{0,1\}^n$.
For $n \ge 1$ we write $b = b_1 + 2b'$, where $b' = \lfloor b/2 \rfloor$ has bit vector $(b_2,\dots,b_n)$.
For $b, c \in \iv{0}{N-1}$, $c \preceq b$ means $c_k \le b_k$ for every $k$.

For a commutative ring $A$, points $x \in A^j$ and $y \in A^{n-j}$, we write $(x,y) \in A^n$ for their concatenation, and $r_{\le j} = (r_1,\dots,r_j)$ for $r \in A^m$ and $j \le m$.
For $x \in A^j$ and $b \in \iv{0}{2^{n-j}-1}$, $(x,b)$ is the point $(x_1,\dots,x_j,b_1,\dots,b_{n-j})$.

\subsection{Fields and univariate polynomials}\label{sec:prelim:fields}

Throughout, $\Fq$ is a finite field of \emph{odd characteristic}, so that $2$ is invertible.
In \cref{sec:sound} the results are also applied over an extension field $\K \supseteq \Fq$, from which the challenges are drawn.
For a field $\F$, $\F[X]_{<d}$ is the space of polynomials of degree less than $d$, and $\coef{a}{P}$ is the coefficient of $X^a$ in $P$.

\begin{lemma}[Root counting, \KB{Lemma~2.1}]\label{lem:roots}
A nonzero $P \in \F[X]$ of degree $d$ has at most $d$ roots in $\F$; two polynomials of $\F[X]_{<d}$ that agree on $d$ distinct points are equal.
\end{lemma}

\begin{lemma}[Even--odd split, \KB{Lemma~2.3}]\label{lem:splits}
Let $d \ge 1$. Every $P \in \F[X]_{<2d}$ can be written uniquely as $P(X) = P_e(X^2) + X P_o(X^2)$ with $P_e, P_o \in \F[X]_{<d}$, namely $\coef{i}{P_e} = \coef{2i}{P}$ and $\coef{i}{P_o} = \coef{2i+1}{P}$.
The map $P \mapsto (P_e, P_o)$ is $\F$-linear.
\end{lemma}

The even--odd split is also defined for every $P \in \F[X]$, by the same coefficient formulas; then $P(X) = P_e(X^2) + X P_o(X^2)$, and $P_e = P_o = 0$ if and only if $P = 0$.

\subsection{Multilinear polynomials}\label{sec:prelim:ml}

Let $A$ be a commutative ring.
A polynomial in $A[x_1,\dots,x_n]$ is \emph{multilinear} if its degree in each variable is at most~$1$; these polynomials form an $A$-module $\mathcal{L}_n(A)$.
A \emph{table} over $A$ with $n$ variables is a function $f : \iv{0}{N-1} \to A$.

\begin{definition}[Monomials, equality and Lagrange polynomials, \KB{Definition~2.4}]\label{def:bases}
For $a \in \iv{0}{N-1}$ the \emph{canonical monomial} is $m_a(x) = \prod_{k=1}^n x_k^{a_k}$.
The \emph{equality polynomials} are $\eq_1(y,z) = yz + (1-y)(1-z)$ and $\eq(y,z) = \prod_{k=1}^n \eq_1(y_k,z_k)$, and the \emph{Lagrange polynomials} are $e_b(x) = \eq(b,x)$ for $b \in \iv{0}{N-1}$.
\end{definition}

The canonical monomials form a basis of $\mathcal{L}_n(A)$.
Every $P \in \mathcal{L}_n(A)$ is written uniquely as $P = P_0 + x_1 P_1$ with $P_0, P_1$ multilinear in $x_2,\dots,x_n$.

\begin{proposition}[Multilinear extension, \KB{Proposition~2.6}]\label{prop:mle}
For every table $f$ there is a unique $\tilde f \in \mathcal{L}_n(A)$ with $\tilde f(b) = f(b)$ for all $b \in B_n$, namely $\tilde f = \sum_{b} f(b)\, e_b$.
The map $f \mapsto \tilde f$ is an $A$-linear bijection from tables onto $\mathcal{L}_n(A)$.
\end{proposition}

\begin{lemma}[Table form and coefficient form, \KB{Lemma~2.9}]\label{lem:mobius}
If $\tilde f = \sum_a \alpha_a m_a$, then $f(b) = \sum_{a \preceq b} \alpha_a$ and $\alpha_a = \sum_{b \preceq a} (-1)^{\wt(a)-\wt(b)} f(b)$, where $\wt$ is the Hamming weight.
\end{lemma}

The table $(\alpha_a)_a$ of \cref{lem:mobius} is the \emph{coefficient form} of $\tilde f$ (\KB{Definition~2.8}), and the map $f \mapsto (\alpha_a)_a$ is the \emph{M\"obius transform}.

\begin{definition}[Restriction, \KB{Definition~2.10}]\label{def:restriction}
Let $n \ge 1$, $f$ a table over $A$ with $n$ variables and $r \in A$. The \emph{restriction} of $f$ at $r$ is the table with $n-1$ variables
\[
  \res_r(f)(b') \;=\; (1-r)\, f(2b') + r\, f(2b'+1), \qquad b' \in \iv{0}{N/2-1}.
\]
\end{definition}

\begin{lemma}[Restriction evaluates the first variable, \KB{Lemma~2.11}]\label{lem:restriction}
$\widetilde{\res_r(f)}(y) = \tilde f(r, y)$ as polynomials in $y = (y_1,\dots,y_{n-1})$.
\end{lemma}

\Cref{sec:restriction} develops the counterpart of \cref{def:restriction,lem:restriction} for the coefficient form, which is the object that the classical FRI fold acts on.

\subsection{Smooth domains and Reed--Solomon codes}\label{sec:prelim:rs}

\begin{definition}[Smooth domain, \KB{Definition~2.14}]\label{def:domain}
A \emph{smooth domain} is a subgroup $L \le \Fq^\times$ of order $M = 2^\mu$ with $\mu \in \N$.
For $j \in \iv{0}{\mu}$, $L^{2^j} = \{\xi^{2^j} : \xi \in L\}$.
\end{definition}

\begin{lemma}[Power maps, \KB{Lemma~2.15}]\label{lem:power}
Let $L$ be a smooth domain of order $M = 2^\mu$ and $j \in \iv{0}{\mu}$.
The map $\xi \mapsto \xi^{2^j}$ is a homomorphism from $L$ onto $L^{2^j}$, the smooth domain of order $M/2^j$, and every element of $L^{2^j}$ has exactly $2^j$ preimages.
If $\mu \ge 1$, then $-1 \in L$ and the preimages of $\xi^2$ under squaring are $\xi$ and $-\xi$, which are distinct.
\end{lemma}

\begin{definition}[Fibres, \KB{Definition~2.16}]\label{def:fibre}
For a smooth domain $L$ of order $M \ge 2$, the \emph{fibre} over $\eta \in L^2$ is the set $\{\xi,-\xi\}$ of its square roots in $L$.
The $M/2$ fibres partition $L$.
\end{definition}

\begin{definition}[Reed--Solomon code, \KB{Definition~2.17}]\label{def:rs}
Let $\F \supseteq \Fq$ be a finite field, $L$ a smooth domain of order $M$ and $1 \le d \le M$.
The evaluation map is $\ev_L : \F[X] \to \F^L$, $P \mapsto (P(\xi))_{\xi \in L}$, and $\RS_\F[L,d] = \ev_L(\F[X]_{<d})$, of rate $\rho = d/M$; we write $\RS[L,d]$ when $\F = \Fq$.
Elements of $\F^L$ are \emph{words}.
The relative Hamming distance is $\Delta(w,w') = |\{\xi : w(\xi) \ne w'(\xi)\}|/M$, and $\Delta(w,\Cc) = \min_{c \in \Cc}\Delta(w,c)$ for a nonempty set $\Cc$ of words.
\end{definition}

\begin{lemma}[Hamming distance, \KB{Lemma~2.18}]\label{lem:metric}
$\Delta$ is a metric on $\F^L$. In particular, if $w$ agrees with $w'$ on at least $\alpha M$ points and $w'$ agrees with $w''$ on at least $\beta M$ points, then $w$ agrees with $w''$ on at least $(\alpha + \beta - 1)M$ points.
\end{lemma}

\begin{lemma}[Parameters of Reed--Solomon codes, \KB{Lemma~2.19}]\label{lem:rs-params}
Let $\Cc = \RS_\F[L,d]$ with $|L| = M$.
\begin{enumerate}
  \item $\ev_L : \F[X]_{<d} \to \Cc$ is an $\F$-linear bijection; for $c \in \Cc$, $P_c$ denotes the polynomial with $\ev_L(P_c) = c$.
  \item Two distinct codewords agree on at most $d-1$ points of $L$.
  \item If $\delta \le (1-\rho)/2$, every word is within distance $\delta$ of at most one codeword.
\end{enumerate}
\end{lemma}

\begin{fact}[Efficient unique decoding, \KB{Fact~2.20}]\label{fact:decode}
There is an algorithm that, given $L$, $d$ and $w \in \F^L$, outputs the unique codeword $c \in \RS_\F[L,d]$ with $\Delta(w,c) < \frac12(1-(d-1)/M)$ whenever it exists, with a number of operations in $\F$ polynomial in $M$.
\end{fact}

\begin{lemma}[Decoding over an extension field, \KB{Lemma~2.21}]\label{lem:descent}
Let $\K \supseteq \Fq$ be a finite field and $w \in \Fq^L$.
If $c \in \RS_\K[L,d]$ satisfies $\Delta(w,c) < \frac12(1-(d-1)/M)$, then $P_c \in \Fq[X]_{<d}$.
\end{lemma}

\paragraph{Even--odd split of words.}
Let $L$ be a smooth domain of order $M \ge 2$ and $w \in \F^L$. The words $w_e, w_o \in \F^{L^2}$ are
\begin{equation}\label{eq:even-odd}
  w_e(\xi^2) \;=\; \frac{w(\xi) + w(-\xi)}{2},
  \qquad
  w_o(\xi^2) \;=\; \frac{w(\xi) - w(-\xi)}{2\xi}
  \qquad (\xi \in L).
\end{equation}

\begin{lemma}[Even--odd split of words, \KB{Lemma~2.22}]\label{lem:split-words}
Let $w \in \F^L$ with $|L| = M \ge 2$.
\begin{enumerate}
  \item $w_e$ and $w_o$ are well defined, and $w \mapsto (w_e,w_o)$ is $\F$-linear.
  \item For $\eta \in L^2$, $w_e(\eta)$ and $w_o(\eta)$ depend only on the values of $w$ on the fibre over $\eta$.
  \item $w(\pm\xi) = w_e(\xi^2) \pm \xi\, w_o(\xi^2)$ for every $\xi \in L$. Consequently, $w$ vanishes on the fibre over $\eta$ if and only if $w_e(\eta) = w_o(\eta) = 0$.
  \item If $1 \le d \le M/2$, $P \in \F[X]_{<2d}$ and $w = \ev_L(P)$, then $w_e = \ev_{L^2}(P_e)$ and $w_o = \ev_{L^2}(P_o)$.
\end{enumerate}
\end{lemma}

\subsection{Correlated agreement}\label{sec:prelim:gaps}

\begin{theorem}[{\cite[Theorem~4.1]{BCIKS23}}; \KB{Theorem~2.23}]\label{thm:bciks}
Let $\F \supseteq \Fq$ be a finite field, $\Cc = \RS_\F[L,d]$ with $|L| = M$ and rate $\rho = d/M$, and $0 < \delta \le (1-\rho)/2$.
Let $u^0, u^1 \in \F^L$ and $S = \{r \in \F : \Delta(u^0 + r u^1, \Cc) \le \delta\}$.
If $|S| > M$, there are a set $L' \subseteq L$ with $|L'| \ge (1-\delta)M$ and codewords $\hat c^0, \hat c^1 \in \Cc$ such that $u^0$ agrees with $\hat c^0$, and $u^1$ with $\hat c^1$, on $L'$.
\end{theorem}

This theorem is the only axiom of the Lean formalisation of KBFold, and of the present one (\cref{app:lean}).

\section{Interactive Oracle Proofs and Polynomial Commitments}\label{sec:proofs}

This section recalls the proof-system notions of KBFold~\cite[\S3]{Mkh26} that we use, and adapts the definition of a polynomial commitment scheme to coefficient-form evaluations and to hypercube sums.
The lemmas are proved in KBFold, and \cref{thm:cdhz} is quoted from Chiesa, Di, Hu and Zheng~\cite{CDHZ25}.

\subsection{Interactive oracle proofs and the probability space}\label{sec:proofs:iop}

A \emph{relation} $\mathcal{R}$ is a set of pairs $(\mathtt{x},\mathtt{w})$ of an instance and a witness, and $L_{\mathcal{R}} = \{\mathtt{x} : \exists\,\mathtt{w},\ (\mathtt{x},\mathtt{w}) \in \mathcal{R}\}$ is its language.
An instance may have an \emph{oracle part}, a word to which the verifier has query access only.

\begin{definition}[Interactive oracle proof, \KB{Definition~3.1}]\label{def:iop}
A \emph{public-coin interactive oracle proof} (IOP)~\cite{BCS16} with $\nu \ge 1$ rounds and finite nonempty challenge sets $\mathsf{Ch}_1,\dots,\mathsf{Ch}_\nu$ is a pair $(\mathsf{P},\mathsf{V})$.
In round $i = 1,\dots,\nu$, $\mathsf{P}$ sends a message $\mathsf{m}_i$, parts of which may be designated as oracles, and $\mathsf{V}$ replies with a challenge $\mathsf{c}_i \in \mathsf{Ch}_i$.
After the last round, $\mathsf{V}$ reads the non-oracle parts of the instance and of the messages, queries the oracles at positions determined by these and by the challenges, and accepts or rejects as a deterministic function of what it has read.
\end{definition}

A \emph{partial transcript} of $i$ rounds is $\tau = (\mathtt{x}, \mathsf{m}_1,\mathsf{c}_1,\dots,\mathsf{m}_i,\mathsf{c}_i)$; for $i = 0$ it is the empty transcript $(\mathtt{x})$, and for $i = \nu$ a full transcript.
Checks made ``during'' a round are predicates on the full transcript, evaluated at the end.
A \emph{deterministic prover} maps each partial transcript of $i-1$ rounds to a message $\mathsf{m}_i$.
Probabilities are over the uniform distribution on $\Omega = \mathsf{Ch}_1 \times \dots \times \mathsf{Ch}_\nu$; $\tau_i$ denotes the partial transcript after $i$ rounds and $\mathsf{m}_i$ the $i$-th message, functions of $(\mathsf{c}_1,\dots,\mathsf{c}_i)$ and $(\mathsf{c}_1,\dots,\mathsf{c}_{i-1})$.

\begin{lemma}[Sequential challenges, \KB{Lemma~3.2}]\label{lem:sequential}
For $i \in \iv{1}{\nu}$ let $E_i \subseteq \mathsf{Ch}_1\times\dots\times\mathsf{Ch}_i$, and suppose that for every $(\mathsf{c}_1,\dots,\mathsf{c}_{i-1})$ at most $k_i$ values $\mathsf{c}_i$ satisfy $(\mathsf{c}_1,\dots,\mathsf{c}_i) \in E_i$.
Then $\Pr_\Omega[\exists\, i : (\mathsf{c}_1,\dots,\mathsf{c}_i) \in E_i] \le \sum_{i} k_i/|\mathsf{Ch}_i|$.
\end{lemma}

\begin{lemma}[Independent queries, \KB{Lemma~3.3}]\label{lem:queries}
Let $\Omega'$ and $L$ be finite nonempty sets, $\kappa \ge 1$, $\varpi \mapsto \mathcal{S}(\varpi) \subseteq L$ and $\mathcal{A} \subseteq \Omega'$.
If $|\mathcal{S}(\varpi)| \le \pi |L|$ for every $\varpi \in \mathcal{A}$, then, under the uniform distribution on $\Omega' \times L^{\times\kappa}$,
$\Pr[\varpi \in \mathcal{A} \text{ and } \xi^{(1)},\dots,\xi^{(\kappa)} \in \mathcal{S}(\varpi)] \le \Pr[\mathcal{A}]\,\pi^\kappa$.
\end{lemma}

\begin{lemma}[Randomised provers, \KB{Lemma~3.4}]\label{lem:randomised}
A bound on the probability of a set of full transcripts that holds for every deterministic prover also holds for every randomised prover whose randomness is independent of the challenges.
\end{lemma}

\begin{definition}[Completeness and soundness, \KB{Definition~3.5}]\label{def:iop-sound}
The IOP is \emph{complete} for $\mathcal{R}$ if the honest prover makes $\mathsf{V}$ accept with probability $1$ on every $(\mathtt{x},\mathtt{w}) \in \mathcal{R}$.
It has \emph{soundness error} $\varepsilon$ if, for every $\mathtt{x} \notin L_{\mathcal{R}}$ and every deterministic prover, $\mathsf{V}$ accepts with probability at most $\varepsilon$.
\end{definition}

\subsection{Round-by-round soundness}\label{sec:proofs:rbr}

\begin{definition}[Round-by-round soundness, \KB{Definition~3.6}]\label{def:rbr}
An IOP with $\nu$ rounds has \emph{round-by-round soundness error} $(\varepsilon_1,\dots,\varepsilon_\nu)$ for $\mathcal{R}$ if there is a set $\Dd$ of partial transcripts, the \emph{doomed} ones, such that:
(1) if $\mathtt{x} \notin L_{\mathcal{R}}$, then $(\mathtt{x}) \in \Dd$;
(2) for every $i$, every doomed $\tau$ of $i-1$ rounds and every message $\mathsf{m}_i$, $\Pr_{\mathsf{c}_i}[(\tau,\mathsf{m}_i,\mathsf{c}_i) \notin \Dd] \le \varepsilon_i$;
(3) $\mathsf{V}$ rejects every doomed full transcript.
\end{definition}

\begin{definition}[Round-by-round knowledge soundness, \KB{Definition~3.7}]\label{def:rbr-knowledge}
An IOP with $\nu$ rounds is \emph{round-by-round knowledge sound} for $\mathcal{R}$ with error $(\varepsilon_1,\dots,\varepsilon_\nu)$ if there are a set $\Dd$ of partial transcripts and a polynomial-time extractor $\mathsf{Ext}$, defined on pairs of a partial transcript and a next message, such that:
(1) $(\mathtt{x}) \in \Dd$ for every instance $\mathtt{x}$;
(2) for every $i$, every doomed $\tau$ of $i-1$ rounds and every $\mathsf{m}_i$, if $\Pr_{\mathsf{c}_i}[(\tau,\mathsf{m}_i,\mathsf{c}_i) \notin \Dd] > \varepsilon_i$, then $(\mathtt{x}, \mathsf{Ext}(\tau,\mathsf{m}_i)) \in \mathcal{R}$;
(3) $\mathsf{V}$ rejects every doomed full transcript.
\end{definition}

\begin{lemma}[Round-by-round implies overall, \KB{Lemma~3.8}]\label{lem:rbr-overall}
Let $\mathsf{P}^*$ be a deterministic prover and $\mathtt{x}$ an instance.
\begin{enumerate}
  \item Under \cref{def:rbr}, if $\mathtt{x} \notin L_{\mathcal{R}}$, then $\mathsf{V}$ accepts with probability at most $\sum_i \varepsilon_i$.
  \item Under \cref{def:rbr-knowledge}, if $\mathsf{V}$ accepts with probability greater than $\sum_i \varepsilon_i$, then there are $i$ and challenges $\mathsf{c}_1,\dots,\mathsf{c}_{i-1}$ with $\tau_{i-1} \in \Dd$ and $(\mathtt{x}, \mathsf{Ext}(\tau_{i-1},\mathsf{m}_i)) \in \mathcal{R}$.
\end{enumerate}
\end{lemma}

\subsection{Coefficient-form commitments and hypercube sums}\label{sec:proofs:pcs}

KBFold commits to a table $f$ and proves claims $\tilde f(z) = v$ about its multilinear extension (\KB{Definition~3.10}).
The scheme of this paper proves claims about the multilinear polynomial whose \emph{coefficients} are the entries of $f$.

\begin{definition}[Coefficient polynomial]\label{def:coeff-poly}
For a table $f$ over a commutative ring $A$ with $n$ variables, the \emph{coefficient polynomial} of $f$ is
\[
  P_f \;=\; \sum_{a=0}^{N-1} f(a)\, m_a \;\in\; \mathcal{L}_n(A).
\]
\end{definition}

By \cref{lem:mobius}, $P_f = \tilde g$ where $g(b) = \sum_{a \preceq b} f(a)$; equivalently, $P_f = \tilde f$ after the M\"obius transform of $f$.
For every $a \in \iv{0}{N-1}$, $m_a(1,\dots,1) = 1$, so
\begin{equation}\label{eq:sum-is-eval}
  \sum_{a=0}^{N-1} f(a) \;=\; P_f(1,\dots,1).
\end{equation}

\begin{definition}[Coefficient-form commitment in the IOP model]\label{def:pcs}
A \emph{coefficient-form commitment scheme} in the IOP model consists of an encoding $\Enc$, mapping each table $f$ to a word $\Enc(f)$, and an \emph{evaluation protocol}, a public-coin IOP whose instance is $(w; z, v)$ with a word $w$ as oracle part, a point $z \in \F^n$ and a value $v \in \F$, for a finite field $\F \supseteq \Fq$.
It is \emph{complete} if the evaluation protocol accepts with probability~$1$ when the prover is honest, $w = \Enc(f)$ and $P_f(z) = v$.
The \emph{sum protocol} is the evaluation protocol on the instances $(w; \mathbf 1, v)$, where $\mathbf 1 = (1,\dots,1)$; by \eqref{eq:sum-is-eval}, it proves $\sum_a f(a) = v$.
\end{definition}

As in KBFold, fix a function $\mathrm{dist}$ on pairs of words and a threshold $\delta$ such that
\begin{equation}\label{eq:dist-unique}
  \text{for every word } w, \text{ at most one table } f \text{ satisfies } \mathrm{dist}(w, \Enc(f)) \le \delta,
\end{equation}
and define the relations
\begin{align*}
  \mathcal{R}^{\delta}_{\mathrm{eval}} &\;=\; \bigl\{\, ((w; z, v), f) \;:\; \mathrm{dist}(w,\Enc(f)) \le \delta \ \text{ and } \ P_f(z) = v \,\bigr\}, \\
  \mathcal{R}^{\delta}_{\mathrm{sum}} &\;=\; \bigl\{\, ((w; v), f) \;:\; \mathrm{dist}(w,\Enc(f)) \le \delta \ \text{ and } \ \textstyle\sum_a f(a) = v \,\bigr\}.
\end{align*}
\Cref{sec:sound} takes for $\mathrm{dist}$ the fibre distance to the encoding, as in KBFold.
By \eqref{eq:sum-is-eval}, $((w;v),f) \in \mathcal{R}^\delta_{\mathrm{sum}}$ if and only if $((w;\mathbf 1, v),f) \in \mathcal{R}^\delta_{\mathrm{eval}}$, so every result for the evaluation protocol specialises to the sum protocol.

\begin{definition}[Evaluation binding, \KB{Definition~3.11}]\label{def:binding}
The evaluation protocol is \emph{$\varepsilon$-evaluation binding} if for every word $w$, point $z$ and values $v \ne v'$, it is not the case that deterministic provers make the verifier accept $(w;z,v)$ and $(w;z,v')$, each with probability greater than $\varepsilon$.
\end{definition}

\begin{lemma}[Knowledge soundness implies binding, \KB{Lemma~3.12}]\label{lem:rbr-binding}
If the evaluation protocol is round-by-round knowledge sound for $\mathcal{R}^\delta_{\mathrm{eval}}$ with error $(\varepsilon_i)_i$ and \eqref{eq:dist-unique} holds, then it is $\varepsilon$-evaluation binding with $\varepsilon = \sum_i \varepsilon_i$.
\end{lemma}

The proof in KBFold uses only \eqref{eq:dist-unique} and \cref{lem:rbr-overall}, and applies verbatim with $P_f$ in place of $\tilde f$.

\subsection{Non-interactive arguments}\label{sec:proofs:bcs}

This subsection concerns only the post-quantum corollary (\cref{cor:pq}).
We use the definitions of KBFold~\cite[\S3.4]{Mkh26}, which restate those of~\cite{CDHZ25}: Merkle trees and their binding (\KB{Lemma~3.13}), relations with an implicit instance and their relaxations $\mathcal{R}_{\le\delta}$, IOPs viewed as interactive oracle reductions (IORs) to the accepting relation $\mathcal{R}_{\mathrm{acc}}$, salted Merkle commitments, committed relations $\mathsf{CR}[\mathcal{R}]$ and relaxed committed relations $\widehat{\mathsf{CR}}[\mathcal{R},\delta]$ (\KB{Definition~3.15}), and the compiler $\mathsf{BCS}[\Pi]$ of~\cite[Construction~11.7]{CDHZ25}.
We restate the hypothesis and the conclusion of the theorem, which are what \cref{sec:sound} checks and uses.

\begin{definition}[Relaxed round-by-round knowledge soundness~\cite{CDHZ25}, \KB{Definition~3.14}]\label{def:rrbr}
A \emph{knowledge state function} for an IOR from $\mathcal{R}$ to $\mathcal{R}_{\mathrm{acc}}$ is a deterministic function $\mathsf{KState}$ that, on a proximity bound $\delta$, a statement $(\mathtt{x}, y)$, a transcript $\mathrm{tr}$ and a witness $\mathtt{w}$, outputs a bit $\mathsf{KState}_\delta(\mathtt{x},y,\mathrm{tr},\mathtt{w})$ such that:
(1) on the empty transcript it is $1$ if and only if $(\mathtt{x},y,\mathtt{w}) \in \mathcal{R}_{\le\delta}$;
(2) if it is $0$ when the prover is about to move, it stays $0$ after every prover message;
(3) on a full transcript it is $1$ if and only if the verifier accepts.
The IOR has \emph{relaxed round-by-round knowledge soundness errors} $(\varepsilon^{\mathrm{rr}}_i)_i$ if there are such a function and a deterministic polynomial-time algorithm $\mathsf{E}$, independent of $\delta$, such that for every $\delta$, statement $(\mathtt{x},y)$, round $i$ and transcript $\mathrm{tr} = (\Pi_1,\mathsf{vr}_1,\dots,\Pi_{i-1},\mathsf{vr}_{i-1},\Pi_i)$,
\[
  \Pr_{\mathsf{vr}_i}\bigl[\, \exists\,\mathtt{w} : \mathsf{KState}_\delta(\mathtt{x},y,\mathrm{tr},\mathsf{E}(\mathtt{x},y,\mathrm{tr}\|\mathsf{vr}_i,\mathtt{w})) = 0 \text{ and } \mathsf{KState}_\delta(\mathtt{x},y,\mathrm{tr}\|\mathsf{vr}_i,\mathtt{w}) = 1 \,\bigr] \;\le\; \varepsilon^{\mathrm{rr}}_i(\mathtt{x},y,\delta).
\]
\end{definition}

\begin{theorem}[{\cite[Theorems~6.10, 8.3 and~11.3]{CDHZ25}}; \KB{Theorem~3.16}]\label{thm:cdhz}
Let $\Pi$ be an IOR from $\mathcal{R}$ to $\mathcal{R}_{\mathrm{acc}}$ with $\nu$ rounds, challenge lengths $\beta_i$ with $\beta_{\min} = \min_i \beta_i$, and relaxed round-by-round knowledge soundness errors $(\varepsilon^{\mathrm{rr}}_i)_i$; let $\varepsilon_{\mathrm{rr}}(\delta)$ be their maximum over the statements under consideration.
Let $l_{\max}$ be the largest length of $y$ and of the messages, $q_{\mathsf{V}}$ the number of random-oracle queries of the verifier of $\mathsf{BCS}[\Pi]$, $q_{\widehat{\mathsf{CR}}}$ that of the decider of $\widehat{\mathsf{CR}}[\mathcal{R},\delta]$, $q_s$ the number of positions read by the IOR verifier, and $T_1 = Q + 1 + q_{\mathsf{V}} + \nu + q_{\widehat{\mathsf{CR}}}$.
Then there is a polynomial-time quantum straightline extractor such that, for every $\delta$ and every query bound $Q$, $\mathsf{BCS}[\Pi]$ is post-quantum straightline knowledge sound from $(\mathsf{CR}[\mathcal{R}], \widehat{\mathsf{CR}}[\mathcal{R},\delta])$, with extraction error
\begin{align*}
  \varepsilon_{\mathrm{ext}}(Q) \;\le\;\;& 320\,(Q+\nu+1)(Q+\nu)\,\varepsilon_{\mathrm{rr}}(\delta) \;+\; \frac{4\,(Q+\nu+1)\,\nu}{2^{\beta_{\min}}} \\
  &+\; \frac{32\, l_{\max} + 1280\, Q\,(2Q+1)^2 + 128\, q_s \lceil \log_2 l_{\max} \rceil}{2^{\lambda_{\mathsf{H}}}} \\
  &+\; \frac{960\,(Q+1)\,T_1^2 + 480\, q_{\mathsf{V}}^2\,(Q + \nu + q_{\mathsf{V}} + 2)}{2^{\lambda_{\mathsf{H}}}},
\end{align*}
and commutativity error $\varepsilon_{\mathrm{com}}(q) \le 240\,(q+1)/2^{\lambda_{\mathsf{H}}}$, where $\lambda_{\mathsf{H}}$ is the output length of the random oracle.
\end{theorem}

The derivation of this specialisation from~\cite{CDHZ25} is given in KBFold; it depends only on the structure of the IOR (number of rounds, message lengths, positions read), not on the encoding.

\section{The Restriction Dictionary}\label{sec:restriction}

This section is the coefficient-form counterpart of the fold dictionary of KBFold (\KB{Theorem~5.3}).
The table is committed as the coefficient vector of a univariate polynomial; the classical FRI fold then restricts one variable of the coefficient polynomial $P_f$ (\cref{def:coeff-poly}), and the scalar claim of the protocol is reduced by the same restriction.
The mechanism is the one recorded in \KB{Remark~5.14} and used by BaseFold~\cite{ZCF24} and DeepFold~\cite{GLHQTZ25}; we state it in the form needed by \cref{sec:scheme,sec:sound}.

Throughout, $A$ is a commutative ring in which $2$ need not be invertible, except where a smooth domain appears.

\subsection{Encoding a table as a coefficient vector}\label{sec:restriction:enc}

\begin{definition}[Coefficient encoding polynomial]\label{def:vf}
For a table $f$ over $A$ with $n$ variables, put
\[
  V_f(X) \;=\; \sum_{a=0}^{N-1} f(a)\, X^{a} \;\in\; A[X]_{<N}.
\]
\end{definition}

The map $f \mapsto V_f$ is an $A$-linear bijection from the tables with $n$ variables onto $A[X]_{<N}$, since the monomials $X^a$, $a \in \iv{0}{N-1}$, form a basis.

\begin{lemma}[Kronecker substitution]\label{lem:kronecker}
For every table $f$ over $A$ with $n$ variables and every $\zeta \in A$,
\[
  V_f(\zeta) \;=\; P_f\bigl(\zeta, \zeta^{2}, \zeta^{4}, \dots, \zeta^{2^{n-1}}\bigr).
\]
In particular $V_f(1) = P_f(1,\dots,1) = \sum_a f(a)$.
\end{lemma}

\begin{proof}
For $a \in \iv{0}{N-1}$, $m_a(\zeta,\zeta^2,\dots,\zeta^{2^{n-1}}) = \prod_{k=1}^n \zeta^{2^{k-1} a_k} = \zeta^{\sum_k a_k 2^{k-1}} = \zeta^a$.
Multiply by $f(a)$ and sum over $a$. The last statement is the case $\zeta = 1$ together with \eqref{eq:sum-is-eval}.
\end{proof}

\subsection{Coefficient restriction}\label{sec:restriction:cres}

\begin{definition}[Coefficient restriction]\label{def:cres}
Let $n \ge 1$, $f$ a table over $A$ with $n$ variables and $\theta \in A$.
The \emph{coefficient restriction} of $f$ at $\theta$ is the table with $n-1$ variables
\[
  \mathrm{cres}_\theta(f)(b') \;=\; f(2b') + \theta\, f(2b'+1), \qquad b' \in \iv{0}{N/2-1}.
\]
For $j \in \iv{0}{n}$ and $\theta = (\theta_1,\dots,\theta_j) \in A^j$, the iterated coefficient restriction is $\mathrm{cres}_{()}(f) = f$ and $\mathrm{cres}_{(\theta_1,\dots,\theta_j)}(f) = \mathrm{cres}_{\theta_j}(\mathrm{cres}_{(\theta_1,\dots,\theta_{j-1})}(f))$, a table with $n-j$ variables.
\end{definition}

The coefficient restriction is $A$-linear in $f$.
It differs from the restriction of \cref{def:restriction}, which reads $(1-r) f(2b') + r f(2b'+1)$: the latter restricts the multilinear extension $\tilde f$, the former restricts the coefficient polynomial $P_f$.

\begin{lemma}[Coefficient restriction evaluates the first variables]\label{lem:cres}
Let $f$ be a table over $A$ with $n$ variables.
\begin{enumerate}
  \item If $n \ge 1$ and $\theta \in A$, then $P_{\mathrm{cres}_\theta(f)}(y) = P_f(\theta, y)$ as polynomials in $y = (y_1,\dots,y_{n-1})$.
  \item For $j \in \iv{0}{n}$ and $\theta \in A^j$, $P_{\mathrm{cres}_\theta(f)}(y) = P_f(\theta, y)$ as polynomials in $y = (y_1,\dots,y_{n-j})$. For $j = n$, the table $\mathrm{cres}_\theta(f)$ has the single entry $P_f(\theta)$.
\end{enumerate}
\end{lemma}

\begin{proof}
(1) Write $a = a_1 + 2a'$. Then $m_a(x_1, y) = x_1^{a_1}\, m_{a'}(y)$, where $m_{a'}$ has $n-1$ variables. Substituting $x_1 = \theta$ and grouping the terms $a_1 = 0$ and $a_1 = 1$,
\[
  P_f(\theta, y) \;=\; \sum_{a'=0}^{N/2-1} \bigl( f(2a') + \theta f(2a'+1) \bigr)\, m_{a'}(y) \;=\; P_{\mathrm{cres}_\theta(f)}(y).
\]
(2) Induction on $j$, as in \KB{Lemma~2.11}: by (1) applied to $\mathrm{cres}_{\theta_{\le j-1}}(f)$ and the induction hypothesis, $P_{\mathrm{cres}_\theta(f)}(y) = P_{\mathrm{cres}_{\theta_{\le j-1}}(f)}(\theta_j, y) = P_f(\theta_{\le j-1}, \theta_j, y)$.
For $j = n$ the polynomial $P$ of a table with $0$ variables is its single entry.
\end{proof}

\subsection{The classical fold}\label{sec:restriction:fold}

\begin{definition}[Classical FRI fold, \KB{Definition~5.10}]\label{def:cfold}
For $\theta \in A$ and $U \in A[X]$, $\cfold_\theta(U) = U_e + \theta\, U_o$, where $U = U_e(X^2) + X U_o(X^2)$ is the even--odd split.
For a word $w$ on a smooth domain $L$ of order at least $2$ over a field $\F \supseteq \Fq$, $\cfold_\theta(w) = w_e + \theta\, w_o \in \F^{L^2}$, with $w_e, w_o$ as in \eqref{eq:even-odd}.
For $\theta = (\theta_1,\dots,\theta_j)$, the iterated fold is $\cfold_{()} = \mathrm{id}$ and $\cfold_{(\theta_1,\dots,\theta_j)} = \cfold_{\theta_j} \circ \cfold_{(\theta_1,\dots,\theta_{j-1})}$, on polynomials and on words.
\end{definition}

If $U \in A[X]_{<2d}$ then $\cfold_\theta(U) \in A[X]_{<d}$, by \cref{lem:splits}.

\begin{theorem}[Restriction dictionary]\label{thm:dictionary}
Let $n \ge 1$, let $f$ be a table over $A$ with $n$ variables, and $\theta \in A$. Then
\[
  \cfold_\theta(V_f) \;=\; V_{\mathrm{cres}_\theta(f)}.
\]
Consequently, for $j \in \iv{0}{n}$ and $\theta \in A^j$, $\cfold_\theta(V_f) = V_{\mathrm{cres}_\theta(f)}$, and for $j = n$ this is the constant polynomial $P_f(\theta_1,\dots,\theta_n)$.
\end{theorem}

\begin{proof}
By \cref{lem:splits}, $\coef{b'}{(V_f)_e} = \coef{2b'}{V_f} = f(2b')$ and $\coef{b'}{(V_f)_o} = f(2b'+1)$ for $b' \in \iv{0}{N/2-1}$, so the coefficient of $X^{b'}$ in $(V_f)_e + \theta (V_f)_o$ is $f(2b') + \theta f(2b'+1) = \mathrm{cres}_\theta(f)(b')$.
The iterated statement follows by induction on $j$, and the case $j = n$ from \cref{lem:cres}(2), since $V_g$ is the constant $g(0)$ for a table $g$ with $0$ variables.
\end{proof}

\begin{figure}[t]
\centering
\begin{tikzcd}[column sep=6em, row sep=2.8em]
  f \in A^{\iv{0}{N-1}} \arrow[r, "\mathrm{cres}_\theta"] \arrow[d, "V"'] &
  \mathrm{cres}_\theta(f) \in A^{\iv{0}{N/2-1}} \arrow[d, "V"] \\
  V_f \in \F[X]_{<N} \arrow[r, "\cfold_\theta"] \arrow[d, "\ev_L"'] &
  \cfold_\theta(V_f) \in \F[X]_{<N/2} \arrow[d, "\ev_{L^2}"] \\
  w \in \F^{L} \arrow[r, "\cfold_\theta"] &
  \cfold_\theta(w) \in \F^{L^2}
\end{tikzcd}
\caption{The restriction dictionary, with $A = \F$. Top square: in the monomial basis the classical fold is the coefficient restriction, which restricts the first variable of $P_f$ (\cref{thm:dictionary,lem:cres}). Bottom square: the word fold is the polynomial fold evaluated on $L^2$ (\cref{lem:cword}). Iterating $n$ times sends the codeword of $f$ to the constant $P_f(\theta_1,\dots,\theta_n)$ (\cref{cor:cword-chain}).}
\label{fig:dictionary}
\end{figure}

\begin{remark}[Relation to the kernel fold]\label{rem:kernel-relation}
KBFold encodes the same table as $U_f = \sum_b f(b) K_b$, where $K_b(X) = \prod_{k=1}^n (X^{2^{k-1}} + b_k)$ (\KB{Definition~4.1}), and its kernel fold $\mathrm{pfold}_r$ restricts $\tilde f$ (\KB{Theorem~5.3}).
The two folds are related by $\cfold_\theta = (1 + 2\theta)\,\mathrm{pfold}_r$ with $r = (1+\theta)/(1+2\theta)$ for $\theta \ne -1/2$ (\KB{Proposition~5.11}).
The two encodings put the sum of the table at different places: $\sum_b f(b) = V_f(1)$ by \cref{lem:kronecker}, and $\sum_b f(b) = \coef{N-1}{U_f}$, since each $K_b$ is monic of degree $\sum_k 2^{k-1} = N-1$.
In odd characteristic, also $\sum_b f(b) = 2^n \tilde f(\tfrac12,\dots,\tfrac12)$, since $\eq(b, (\tfrac12,\dots,\tfrac12)) = 2^{-n}$ for every $b \in B_n$; this is the claim that the KBFold scheme proves for a sum.
\end{remark}

\subsection{Folding words}\label{sec:restriction:words}

In this subsection $\F \supseteq \Fq$ is a finite field and $L$ is a smooth domain of order $M \ge 2$, so that \cref{lem:split-words} applies.

\begin{lemma}[The classical word fold]\label{lem:cword}
Let $w \in \F^L$ and $\theta \in \F$.
\begin{enumerate}
  \item (Linearity and locality.) $w \mapsto \cfold_\theta(w)$ is $\F$-linear, and $\cfold_\theta(w)(\eta)$ depends only on the values of $w$ on the fibre over $\eta \in L^2$.
  \item (Line form.) $\cfold_\theta(w) = u^0 + \theta\, u^1$ with $u^0 = w_e$ and $u^1 = w_o$, which do not depend on $\theta$; and $w$ vanishes on the fibre over $\eta$ if and only if $u^0(\eta) = u^1(\eta) = 0$.
  \item (Explicit formula.) For $\xi \in L$,
  \[
    \cfold_\theta(w)(\xi^2) \;=\; \frac{(\xi + \theta)\, w(\xi) + (\xi - \theta)\, w(-\xi)}{2\xi}.
  \]
  \item (Consistency with the polynomial fold.) If $1 \le d \le M/2$, $U \in \F[X]_{<2d}$ and $w = \ev_L(U)$, then $\cfold_\theta(w) = \ev_{L^2}(\cfold_\theta(U)) \in \RS_\F[L^2, d]$.
\end{enumerate}
\end{lemma}

\begin{proof}
(1) and (2) follow from \cref{lem:split-words}(1)--(3).
(3) Substitute \eqref{eq:even-odd}: $\frac{w(\xi)+w(-\xi)}{2} + \theta \frac{w(\xi)-w(-\xi)}{2\xi}$ has the stated numerator over $2\xi$.
(4) By \cref{lem:split-words}(4), $w_e = \ev_{L^2}(U_e)$ and $w_o = \ev_{L^2}(U_o)$; $\ev_{L^2}$ is linear, and $\cfold_\theta(U) \in \F[X]_{<d}$ with $d \le |L^2|$.
\end{proof}

Item 2 is the property through which correlated agreement enters the soundness proof (\cref{sec:sound}).
For the kernel fold, KBFold has to rewrite the fold as a line with generators $w_o - w_e$ and $2w_e - w_o$ (\KB{Lemma~5.7}); for the classical fold the generators are the even and odd parts themselves.

\begin{corollary}[Folding a committed codeword]\label{cor:cword-chain}
Let $L$ be a smooth domain of order $M \ge N = 2^n$, $f$ a table over $\F$ with $n$ variables, $w_0 = \ev_L(V_f)$, and $\theta \in \F^n$.
For $j \in \iv{0}{n}$ put $w_j = \cfold_{\theta_{\le j}}(w_0) \in \F^{L^{2^j}}$. Then
\[
  w_j \;=\; \ev_{L^{2^j}}\bigl(V_{\mathrm{cres}_{\theta_{\le j}}(f)}\bigr) \;\in\; \RS_\F\bigl[L^{2^j}, N/2^j\bigr],
\]
and $w_n$ is the constant word on $L^{2^n}$ with value $P_f(\theta_1,\dots,\theta_n)$.
\end{corollary}

\begin{proof}
Induction on $j$. For the step, apply \cref{lem:cword}(4) on the domain $L^{2^{j-1}}$, of order $M/2^{j-1} \ge 2$ (\cref{lem:power}), with $2d = N/2^{j-1} \le M/2^{j-1}$, and use \cref{thm:dictionary}. The value of $w_n$ is given by \cref{thm:dictionary} with $j = n$.
\end{proof}

\subsection{The round polynomial}\label{sec:restriction:round}

The scalar claim of the protocol has the form $P_g(z) = v$ for the current table $g$ and the remaining coordinates $z$ of the evaluation point.
One round replaces the first coordinate of $z$ by a fresh challenge.

\begin{definition}[Honest round polynomial]\label{def:round-poly}
Let $k \ge 1$, $g$ a table over $A$ with $k$ variables, and $y \in A^{k-1}$.
Let $g_e$ and $g_o$ be the tables with $k-1$ variables $g_e(b') = g(2b')$ and $g_o(b') = g(2b'+1)$.
The \emph{honest round polynomial} of $(g, y)$ is
\[
  h_{g,y}(T) \;=\; P_{g_e}(y) + T\, P_{g_o}(y) \;\in\; A[T]_{<2}.
\]
\end{definition}

\begin{lemma}[Round polynomial]\label{lem:round-poly}
With the notation of \cref{def:round-poly}:
\begin{enumerate}
  \item $h_{g,y}(t) = P_g(t, y)$ for every $t \in A$; in particular $h_{g,y}(t) = P_{\mathrm{cres}_t(g)}(y)$.
  \item For $y = (1,\dots,1)$, $h_{g,y}(T) = U_e(1) + T\, U_o(1)$, where $U = V_g$ and $U = U_e(X^2) + XU_o(X^2)$.
  \item Let $A = \F$ be a field and $s \in \F[T]_{<2}$ with $s \ne h_{g,y}$. Then $s(\theta) = h_{g,y}(\theta)$ for at most one $\theta \in \F$.
\end{enumerate}
\end{lemma}

\begin{proof}
(1) $\mathrm{cres}_t(g) = g_e + t\, g_o$ and $P$ is linear in the table, so $P_{\mathrm{cres}_t(g)}(y) = P_{g_e}(y) + t P_{g_o}(y) = h_{g,y}(t)$; and $P_{\mathrm{cres}_t(g)}(y) = P_g(t,y)$ by \cref{lem:cres}(1).
(2) By \cref{lem:splits}, $\coef{b'}{U_e} = g(2b') = g_e(b')$ and $\coef{b'}{U_o} = g(2b'+1) = g_o(b')$, so $U_e = V_{g_e}$ and $U_o = V_{g_o}$; by \cref{lem:kronecker}, $V_{g_e}(1) = P_{g_e}(\mathbf 1)$ and $V_{g_o}(1) = P_{g_o}(\mathbf 1)$.
(3) $s - h_{g,y}$ is a nonzero polynomial of degree at most $1$; apply \cref{lem:roots}.
\end{proof}

By \cref{lem:round-poly}(1), $h_{g,y}(z_1) = P_g(z_1, y)$, so the claim $P_g(z_1, y) = v$ is checked as $h_{g,y}(z_1) = v$, and the new claim after the challenge $\theta$ is $P_{\mathrm{cres}_\theta(g)}(y) = h_{g,y}(\theta)$.
This is the restriction step of the protocol of \cref{sec:scheme}.
It replaces the round of the sumcheck protocol of KBFold (\KB{Lemma~3.9}), whose round polynomials have degree $2$ because they carry the factor $\eq(\cdot, z)$, by an affine polynomial with no equality table.

\begin{example}\label{ex:f17}
Over $\mathbb{F}_{17}$ with $n = 2$, let $f = (1, 2, 3, 4)$, so $V_f = 1 + 2X + 3X^2 + 4X^3$, $P_f(x_1,x_2) = 1 + 2x_1 + 3x_2 + 4x_1x_2$, and $\sum_a f(a) = 10 = V_f(1)$.
For the sum, $z = (1,1)$ and the first round polynomial is $h(T) = (1 + 3) + T(2 + 4) = 4 + 6T$, with $h(1) = 10$.
With $\theta_1 = 5$: $\mathrm{cres}_5(f) = (1 + 5\cdot 2,\ 3 + 5 \cdot 4) = (11, 6)$, so $\cfold_5(V_f) = 11 + 6X$, and the new claim is $h(5) = 34 = 0 = 11 + 6$.
The second round polynomial is $11 + 6T$; with $\theta_2 = 7$ the final constant is $11 + 42 = 53 = 2$, and indeed $P_f(5, 7) = 1 + 10 + 21 + 140 = 172 = 2$ in $\mathbb{F}_{17}$.
A prover claiming the sum $11$ must send an affine $s$ with $s(1) = 11$, for instance $s(T) = 1 + 10T$; then $s - h = 4T - 3$ vanishes only at $T = 5$, so the false claim survives the first round only if $\theta_1 = 5$ (\cref{lem:round-poly}(3)).
\end{example}

\section{The Commitment Scheme}\label{sec:scheme}

We now build the coefficient-form commitment scheme of \cref{def:pcs} on the restriction dictionary.
The prover commits to a table $f$ as the Reed--Solomon codeword of $V_f$, with no change of basis.
A claim $P_f(z) = v$ is proved by restriction rounds (\cref{lem:round-poly}) whose challenges are also the fold challenges; for $z = \mathbf 1$ this proves $\sum_a f(a) = v$.
The structure, the stopping parameter and the committed levels are those of the KBFold scheme (\KB{\S6}), with the classical fold in place of the kernel fold and the restriction step in place of the sumcheck.

\subsection{Parameters and encoding}\label{sec:scheme:params}

\paragraph{Fields.}
The smooth domain lies in the base field $\Fq$; tables, words, the evaluation point and the challenges live in a finite field $\F \supseteq \Fq$, the \emph{challenge field}.
Taking $\F = \K$, an extension of $\Fq$, gives the scheme with extension-field challenges of \cref{sec:impl}, in which honest tables have entries in $\Fq$.
All results of \cref{sec:scheme,sec:sound} hold for every such $\F$, with $|\F|$ in the error bounds.

\paragraph{Parameters.}
As in \KB{\S6.1}:
\begin{itemize}
  \item the number of variables $n \ge 1$, and $N = 2^n$;
  \item a rate $\rho = 2^{-R}$ with $R \ge 1$, and a smooth domain $L \le \Fq^\times$ of order $M = N/\rho = 2^{n+R}$;
  \item a number of folding rounds $\ell \in \iv{0}{n}$, and a number of queries $\kappa \ge 1$.
\end{itemize}
For $j \in \iv{0}{\ell}$ let $L_j = L^{2^j}$, of order $M_j = M/2^j \ge 2$, let $N_j = N/2^j$, and $\Cc_j = \RS_\F[L_j, N_j]$; every $\Cc_j$ has rate $\rho$, and $L_{j+1} = (L_j)^2$.

\begin{definition}[Encoding]\label{def:encoding}
The \emph{encoding} of a table $f : \iv{0}{N-1} \to \F$ is $\Enc(f) = \ev_L(V_f) \in \Cc_0$, with $V_f$ as in \cref{def:vf}.
\end{definition}

\begin{lemma}[Encoding]\label{lem:encoding}
$\Enc$ is an $\F$-linear bijection from the tables over $\F$ onto $\Cc_0$; for $c \in \Cc_0$, the table $f$ with $\Enc(f) = c$ is the coefficient vector of $P_c$, $f(a) = \coef{a}{P_c}$.
\end{lemma}

\begin{proof}
$f \mapsto V_f$ is a linear bijection onto $\F[X]_{<N}$, and $\ev_L$ is a linear bijection from $\F[X]_{<N}$ onto $\Cc_0$ with inverse $c \mapsto P_c$ (\cref{lem:rs-params}(1)).
\end{proof}

$\Enc$ is computed by one number-theoretic transform (NTT) of size $M$ on the zero-padded table: the table \emph{is} the coefficient vector, and no transform precedes the NTT.
KBFold spends $\tfrac{n}{2}N$ additions on the kernel transform (\KB{Corollary~4.11}), and a coefficient-form scheme proving claims about $\tilde f$ spends as many on the M\"obius transform (\cref{lem:mobius}).

\subsection{The evaluation protocol}\label{sec:scheme:protocol}

The commitment to $f$ is $w_0 = \Enc(f)$.
The evaluation protocol $\Pieval$ proves the claim $P_f(z) = v$ for a public point $z = (z_1,\dots,z_n) \in \F^n$ and a public value $v \in \F$; its instance is $(w_0; z, v)$, with oracle access to $w_0$.
The \emph{sum protocol} $\Pisum$ is $\Pieval$ with $z = \mathbf 1 = (1,\dots,1)$; it proves $\sum_a f(a) = v$ (\cref{def:pcs}).

\paragraph{Round structure.}
$\Pieval$ is a public-coin IOP with $\ell + 1$ rounds.
\begin{itemize}
  \item In round $j \in \iv{1}{\ell}$, the prover sends a \emph{round polynomial} $s_j \in \F[T]_{<2}$, together with, if $j \ge 2$, an oracle $w_{j-1} \in \F^{L_{j-1}}$; the verifier replies with a uniformly random challenge $\theta_j \in \F$.
  \item In round $\ell + 1$, the prover sends a \emph{final table} $g : \iv{0}{N_\ell - 1} \to \F$ in the clear; the verifier replies with $\kappa$ independent uniformly random query points $\xi^{(1)}_0,\dots,\xi^{(\kappa)}_0 \in L$, and decides.
\end{itemize}
The challenge sets are $\mathsf{Ch}_j = \F$ for $j \in \iv{1}{\ell}$ and $\mathsf{Ch}_{\ell+1} = L^{\times\kappa}$.
A round polynomial is transmitted as its two coefficients; \cref{rem:compress} halves this.

\paragraph{Honest prover.}
The honest prover keeps one table, initialised to $f_0 = f$, and after round $j$ replaces it by its coefficient restriction $f_j = \mathrm{cres}_{\theta_j}(f_{j-1})$ (\cref{def:cres}).
In round $j$ it sends the honest round polynomial of $f_{j-1}$ at the remaining coordinates (\cref{def:round-poly}),
\begin{equation}\label{eq:honest-sj}
  s_j(T) \;=\; h_{f_{j-1},\, z_{>j}}(T) \;=\; P_{(f_{j-1})_e}(z_{>j}) + T\, P_{(f_{j-1})_o}(z_{>j}),
  \qquad z_{>j} = (z_{j+1},\dots,z_n),
\end{equation}
and, if $j \ge 2$, the oracle $w_{j-1} = \cfold_{\theta_{j-1}}(w_{j-2})$ (\cref{def:cfold}).
In round $\ell + 1$ it sends $g = f_\ell$.
For $\Pisum$, $s_j(T) = \sum_{b} f_{j-1}(2b) + T \sum_b f_{j-1}(2b+1)$, by \cref{lem:round-poly}(2) and \cref{lem:kronecker}.

\paragraph{Verifier.}
Set $v_0 = v$ and $v_j = s_j(\theta_j)$ for $j \in \iv{1}{\ell}$.
Let $G = V_g \in \F[X]_{<N_\ell}$ and $w_\ell = \ev_{L_\ell}(G)$; the verifier computes values of $w_\ell$ itself.
The verifier accepts if and only if all of the following checks pass.
\begin{enumerate}
  \item \emph{Restriction.} $s_j(z_j) = v_{j-1}$ for every $j \in \iv{1}{\ell}$.
  \item \emph{Closure.} $v_\ell = P_g(z_{\ell+1},\dots,z_n)$.
  \item \emph{Queries.} For every $i \in \iv{1}{\kappa}$, with $\xi_0 = \xi_0^{(i)}$ and $\xi_j = \xi_0^{2^j}$:
  \begin{itemize}
    \item if $\ell \ge 1$: for every $j \in \iv{1}{\ell}$, the value $\cfold_{\theta_j}(w_{j-1})(\xi_j)$, computed by \cref{lem:cword}(3) from the two queried values $w_{j-1}(\pm\xi_{j-1})$, equals $w_j(\xi_j)$ (a query to the oracle $w_j$ if $j < \ell$, and $G(\xi_\ell)$ if $j = \ell$);
    \item if $\ell = 0$: $w_0(\xi_0) = G(\xi_0)$.
  \end{itemize}
\end{enumerate}
For $\Pisum$, the restriction check reads $s_j(1) = v_{j-1}$ and the closure check reads $v_\ell = \sum_b g(b)$.
Unlike in KBFold (\KB{\S6.2}), the closure check carries no factor $\prod_j \eq_1(\theta_j, z_j)$: the reduced claim after round $j$ is directly a claim on $P_{f_j}$.

\begin{remark}[Stopping parameter]\label{rem:stop}
With $\ell = n$, the final table is the single value $P_f(\theta_1,\dots,\theta_n)$ and $G$ is a constant; with $\ell = 0$, the prover sends $f$ in the clear.
Intermediate values trade the $N_\ell = 2^{n-\ell}$ field elements of the final message against $\ell$ folding rounds, as in \KB{Remark~6.3}.
\end{remark}

\begin{remark}[One field element per round]\label{rem:compress}
Since the verifier checks $s_j(z_j) = v_{j-1}$, the prover may send only the linear coefficient $b_j$ of $s_j(T) = a_j + b_j T$, the verifier setting $a_j = v_{j-1} - z_j b_j$; the restriction check then holds by construction, and the protocol is the same IOP up to this bijective re-encoding of the messages that pass the check.
Every soundness result below therefore holds for the compressed protocol.
\end{remark}

\subsection{Committing to fewer words}\label{sec:scheme:skip}

As in \KB{\S6.3}, the prover may commit only to the words $w_j$ for $j$ in a public set $\mathcal{J} \subseteq \iv{0}{\ell-1}$ with $0 \in \mathcal{J}$, and the verifier folds the intermediate levels itself.
For $j \in \mathcal{J}$ let $j^+$ be the next element of $\mathcal{J}$, or $\ell$, let $d_j = j^+ - j$, let $S_j(\xi) = \{\zeta \in L_j : \zeta^{2^{d_j}} = \xi^{2^{d_j}}\}$, and $\cfold^{(j)}_\theta = \cfold_{\theta_{j^+}} \circ \dots \circ \cfold_{\theta_{j+1}}$.

\begin{lemma}[Local folding, \KB{Lemma~6.4}]\label{lem:local-fold}
For $j \in \mathcal{J}$, $w \in \F^{L_j}$ and $\xi \in L_j$, the value $\cfold^{(j)}_\theta(w)(\xi^{2^{d_j}})$ depends only on the values of $w$ on $S_j(\xi)$, and is computed from them by $2^{d_j} - 1$ applications of \cref{lem:cword}(3).
\end{lemma}

The protocol $\Pieval^{\mathcal{J}}$ is $\Pieval$ in which the oracle $w_{j-1}$ is sent only if $j - 1 \in \mathcal{J}$, and in which, for every query point and every $j \in \mathcal{J}$, the verifier reads $w_j$ on $S_j(\xi_j)$, computes $\cfold^{(j)}_\theta(w_j)(\xi_{j^+})$, and compares it with $w_{j^+}(\xi_{j^+})$, or $G(\xi_\ell)$ if $j^+ = \ell$.
For $\mathcal{J} = \iv{0}{\ell-1}$ it is $\Pieval$.

\begin{lemma}[Skipping commitments, \KB{Lemma~6.5}]\label{lem:skip}
For every prover $\tilde{\mathsf{P}}$ of $\Pieval^{\mathcal{J}}$, the prover of $\Pieval$ that sends the messages of $\tilde{\mathsf{P}}$ and, in round $j+1$ for $j \in \iv{1}{\ell-1} \setminus \mathcal{J}$, the oracle $w_j = \cfold_{\theta_j}(w_{j-1})$, is accepted, for every instance and every choice of the challenges, whenever $\tilde{\mathsf{P}}$ is.
\end{lemma}

The proofs of \KB{Lemmas~6.4 and~6.5} use only the locality and the composition of the word fold, and apply verbatim to $\cfold$ by \cref{lem:cword}(1).
\Cref{lem:skip} transfers soundness and evaluation binding from $\Pieval$ to $\Pieval^{\mathcal{J}}$ with the same errors.

\subsection{Completeness}\label{sec:scheme:complete}

\begin{lemma}[Honest prover]\label{lem:honest}
Let $f$ be a table over $\F$, $z \in \F^n$ and $\theta_1,\dots,\theta_\ell \in \F$, and let $f_j$ and $s_j$ be computed by the honest prover. Then, for $j \in \iv{0}{\ell}$:
\begin{enumerate}
  \item $f_j = \mathrm{cres}_{\theta_{\le j}}(f)$, and $P_{f_j}(y) = P_f(\theta_{\le j}, y)$ for every $y \in \F^{n-j}$;
  \item for $j \ge 1$, $s_j \in \F[T]_{<2}$, $s_j(z_j) = P_{f_{j-1}}(z_{\ge j})$ and $s_j(\theta_j) = P_{f_j}(z_{>j})$.
\end{enumerate}
\end{lemma}

\begin{proof}
(1) is the definition and \cref{lem:cres}(2).
(2) By \cref{lem:round-poly}(1), $s_j(t) = P_{f_{j-1}}(t, z_{>j})$ for every $t$; take $t = z_j$, and $t = \theta_j$ together with $P_{f_{j-1}}(\theta_j, y) = P_{\mathrm{cres}_{\theta_j}(f_{j-1})}(y) = P_{f_j}(y)$ (\cref{lem:cres}(1)).
\end{proof}

\begin{theorem}[Completeness]\label{thm:complete}
If $w_0 = \Enc(f)$, $v = P_f(z)$ and the prover is honest, the verifier of $\Pieval$ (and of $\Pieval^{\mathcal{J}}$) accepts with probability $1$.
In particular, if $v = \sum_a f(a)$, the verifier of $\Pisum$ accepts with probability $1$.
\end{theorem}

\begin{proof}
\emph{Restriction.} By \cref{lem:honest}(2) and induction, $v_j = P_{f_j}(z_{>j})$ for every $j \in \iv{0}{\ell}$: for $j = 0$ it is the hypothesis $v = P_f(z)$, and $v_j = s_j(\theta_j) = P_{f_j}(z_{>j})$. Hence $s_j(z_j) = P_{f_{j-1}}(z_{\ge j}) = v_{j-1}$.

\emph{Closure.} $v_\ell = P_{f_\ell}(z_{>\ell}) = P_g(z_{>\ell})$, since $g = f_\ell$.

\emph{Queries.} If $\ell = 0$, $G = V_f$ and $w_0 = \ev_L(G)$.
If $\ell \ge 1$, extend $(\theta_1,\dots,\theta_\ell)$ arbitrarily to $\theta \in \F^n$; by \cref{cor:cword-chain}, $\cfold_{\theta_{\le j}}(w_0) = \ev_{L_j}(V_{f_j})$ for $j \in \iv{0}{\ell}$.
The honest oracles are these words for $j < \ell$, and $V_{f_\ell} = V_g = G$.
Hence $\cfold_{\theta_j}(w_{j-1}) = w_j$ on all of $L_j$ for every $j \in \iv{1}{\ell}$, and every query check passes; for $\Pieval^{\mathcal{J}}$, the same identities give $\cfold^{(j)}_\theta(w_j) = w_{j^+}$.
The last statement is the case $z = \mathbf 1$, by \cref{lem:kronecker}.
\end{proof}

\subsection{Costs}\label{sec:scheme:costs}

Let $t_{\mathsf{H}}$ be the cost of one hash evaluation; field operations are counted up to constant factors.

\begin{itemize}
  \item \emph{Commit.} One NTT of size $M$, $O(M \log M)$ operations, and a Merkle tree over the $M/2$ fibres of $w_0$, $O(M) \cdot t_{\mathsf{H}}$. No basis transform.
  \item \emph{Prover, evaluation.} The tensors $\tau_j = \bigl(m_a(z_{>j})\bigr)_a$, $j = n-1, \dots, 0$, are computed backwards by $\tau_{j-1}(a_1 + 2a') = z_j^{a_1}\, \tau_j(a')$, in $O(N)$ operations in total; round $j$ then costs $O(N_{j-1})$ for $s_j$ (two inner products with $\tau_j$) and for $f_j$, hence $O(N)$ over all rounds.
  For $\Pisum$, $\tau_j = \mathbf 1$ and the round polynomials cost only additions.
  The folds cost $O(M)$ in total, and the Merkle trees of $w_1,\dots,w_{\ell-1}$ cost $O(M)\cdot t_{\mathsf{H}}$.
  \item \emph{Verifier.} Restriction checks: $O(\ell)$. Closure: $O(N_\ell)$ (a sum for $\Pisum$). Queries: $O(\kappa\ell)$ field operations, $O(\kappa\ell\log M)\cdot t_{\mathsf{H}}$ for the authentication paths, and $\kappa$ evaluations of $G$, $O(\kappa N_\ell)$.
  \item \emph{Proof size.} For $\ell \ge 1$: $\ell$ field elements for the round polynomials (\cref{rem:compress}), against $3\ell$ for the degree-$2$ sumcheck of KBFold; $N_\ell$ for the final table; $\ell - 1$ Merkle roots; $\kappa\ell$ opened fibres with their paths. With committed levels $\mathcal{J}$, the openings are those of \KB{\S6.4}.
\end{itemize}

\begin{remark}[Relation to DeepFold, BaseFold and Gemini]\label{rem:deepfold}
For $\mathcal{J} = \iv{0}{\ell-1}$, $\ell = n$ and without its out-of-domain points, the evaluation protocol of DeepFold~\cite{GLHQTZ25} has the round structure of $\Pieval$: the prover sends $P_{f_{j-1}}(T, z_{>j})$, the verifier checks it at $z_j$ and evaluates it at the fold challenge.
BaseFold~\cite{ZCF24} proves a coefficient-form claim with a degree-$2$ sumcheck instead.
Gemini~\cite{BCHO22} folds at the coordinates of $z$ rather than at random challenges, and therefore commits to every folded polynomial and tests each one against the previous; at $z = \mathbf 1$ its folds compute $\sum_a f(a)$ by repeated pairwise sums.
In $\Pieval$ the folds are random and the scalar claim follows them, so the folded words need no proof beyond the query checks of FRI.
\end{remark}

\section{Soundness}\label{sec:sound}

We prove that $\Pieval$ is sound, evaluation binding, round-by-round knowledge sound and, once compiled, post-quantum knowledge sound, in the unique-decoding regime.
The structure is that of \KB{\S7}.
The statements about fibre distance and decoding concern only the codes and carry over unchanged (\cref{sec:sound:fibre}).
The two folding lemmas are proved again for the classical fold (\cref{sec:sound:lemmas}); correlated agreement (\cref{thm:bciks}) enters through the line form of \cref{lem:cword}(2), whose generators are the even and odd parts of the word.
The passing sets and the codeword chain carry over with $\cfold$ in place of the kernel fold (\cref{sec:sound:thm}).
The sumcheck step of KBFold is replaced by a restriction step (\cref{lem:round-poly}(3)), which costs one bad challenge per round instead of two.

Throughout this section we use the notation of \cref{sec:scheme}: the challenge field $\F \supseteq \Fq$, the domains $L_j$ of order $M_j$, and the codes $\Cc_j = \RS_\F[L_j, N_j]$, all of rate $\rho$.
We fix a proximity parameter $0 < \delta \le (1-\rho)/2$.
We assume $1 \le \ell \le n$, except in \cref{sec:sound:l0}, \cref{prop:rbr-l0} and the third item of \cref{rem:fs}, which treat $\ell = 0$.
Probabilities are over the challenge space $\F^\ell \times L^{\times\kappa}$ of $\Pieval$ (\cref{sec:scheme:protocol}).

\subsection{Fibre distance and decoding}\label{sec:sound:fibre}

For $j \in \iv{0}{\ell-1}$, the domain $L_j$ has order $M_j \ge 2$, and its fibres (\cref{def:fibre}) are indexed by the $M_{j+1}$ points of $L_{j+1} = (L_j)^2$.

\begin{definition}[Fibre distance, \KB{Definition~7.1}]\label{def:fibre-dist}
Let $j \in \iv{0}{\ell-1}$. For $w, c \in \F^{L_j}$,
\[
  \Delta^{\mathrm{fib}}_j(w,c) \;=\; \frac{\bigl|\{\, \eta \in L_{j+1} \;:\; w \text{ and } c \text{ differ somewhere on the fibre over } \eta \,\}\bigr|}{M_{j+1}},
\]
and $\Delta^{\mathrm{fib}}_j(w, \Cc_j) = \min_{c \in \Cc_j} \Delta^{\mathrm{fib}}_j(w,c)$.
\end{definition}

\begin{lemma}[Fibre distance, \KB{Lemma~7.2}]\label{lem:fibre-dist}
Let $j \in \iv{0}{\ell-1}$ and $w, c \in \F^{L_j}$.
\begin{enumerate}
  \item $\Delta(w,c) \le \Delta^{\mathrm{fib}}_j(w,c)$.
  \item $\Delta^{\mathrm{fib}}_j(w,c) \le \delta$ if and only if $w$ and $c$ agree on every fibre over some set of at least $(1-\delta)M_{j+1}$ points of $L_{j+1}$.
\end{enumerate}
\end{lemma}

\begin{proof}
(1) Each point of $L_j$ at which $w$ and $c$ differ lies in a fibre over which they differ, and each of these $\Delta^{\mathrm{fib}}_j(w,c)\,M_{j+1}$ fibres has two points, so $\Delta(w,c)\,M_j \le 2\,\Delta^{\mathrm{fib}}_j(w,c)\,M_{j+1} = \Delta^{\mathrm{fib}}_j(w,c)\,M_j$.
(2) The set of points over whose fibre $w$ and $c$ agree has $(1 - \Delta^{\mathrm{fib}}_j(w,c))\,M_{j+1}$ elements.
\end{proof}

\begin{lemma}[Unique fibre decoding, \KB{Lemma~7.3}]\label{lem:unique}
Let $j \in \iv{0}{\ell-1}$. For every $w \in \F^{L_j}$ there is at most one $c \in \Cc_j$ with $\Delta^{\mathrm{fib}}_j(w,c) \le \delta$.
\end{lemma}

\begin{proof}
Such a $c$ satisfies $\Delta(w,c) \le \delta \le (1-\rho)/2$ by \cref{lem:fibre-dist}(1); apply \cref{lem:rs-params}(3).
\end{proof}

\begin{definition}[Decoded codeword and decoded table]\label{def:decoded}
Let $j \in \iv{0}{\ell-1}$ and $w \in \F^{L_j}$ with $\Delta^{\mathrm{fib}}_j(w, \Cc_j) \le \delta$.
The \emph{decoded codeword} $\mathrm{dec}_j(w)$ is the unique $c \in \Cc_j$ of \cref{lem:unique}.
For $j = 0$, the \emph{decoded table} of $w$ is the unique table $f^*$ over $\F$ with $\Enc(f^*) = \mathrm{dec}_0(w)$, that is, the coefficient vector $f^*(a) = \coef{a}{P_{\mathrm{dec}_0(w)}}$ (\cref{lem:encoding}).
\end{definition}

\begin{remark}[Computing the decoded table]\label{rem:ext}
Since $\Delta(w, \mathrm{dec}_0(w)) \le \delta < \tfrac12(1 - (N-1)/M)$ (\cref{lem:rs-params}), the algorithm of \cref{fact:decode} finds $\mathrm{dec}_0(w)$, and an inverse NTT gives $P_{\mathrm{dec}_0(w)}$, whose coefficient vector is $f^*$; unlike in KBFold (\KB{Remark~7.5}), no basis transform follows.
For an arbitrary $w$, if the algorithm of \cref{fact:decode} returns a codeword $c$, checking whether $\Delta^{\mathrm{fib}}_0(w, c) \le \delta$ costs $O(M)$ comparisons; this decides whether $\Delta^{\mathrm{fib}}_0(w, \Cc_0) \le \delta$.
If $w \in \Fq^L$, then $f^*$ has entries in $\Fq$, by \cref{lem:descent}.
\end{remark}

\subsection{Two folding lemmas}\label{sec:sound:lemmas}

\begin{lemma}[Folding far words]\label{lem:far}
Let $j \in \iv{0}{\ell-1}$ and $w \in \F^{L_j}$ with $\Delta^{\mathrm{fib}}_j(w, \Cc_j) > \delta$. Then
\[
  \bigl|\{\, \theta \in \F \;:\; \Delta\bigl(\cfold_\theta(w),\, \Cc_{j+1}\bigr) \le \delta \,\}\bigr| \;\le\; M_{j+1}.
\]
\end{lemma}

\begin{proof}
Suppose the set has more than $M_{j+1}$ elements. By \cref{lem:cword}(2), $\cfold_\theta(w) = u^0 + \theta u^1$ with $u^0 = w_e$ and $u^1 = w_o$.
Apply \cref{thm:bciks} to the code $\Cc_{j+1}$, of length $M_{j+1}$ and rate $\rho$: there are $L' \subseteq L_{j+1}$ with $|L'| \ge (1-\delta)M_{j+1}$ and codewords $\hat c^0, \hat c^1 \in \Cc_{j+1}$ with $w_e = \hat c^0$ and $w_o = \hat c^1$ on $L'$.
Let $T(X) = P_{\hat c^0}(X^2) + X\, P_{\hat c^1}(X^2)$, which lies in $\F[X]_{<N_j}$ since $P_{\hat c^0}, P_{\hat c^1} \in \F[X]_{<N_{j+1}}$ and $N_j = 2N_{j+1}$.
For every $\eta \in L'$ and $\xi \in L_j$ with $\xi^2 = \eta$, \cref{lem:split-words}(3) gives
$w(\pm \xi) = w_e(\eta) \pm \xi\, w_o(\eta) = P_{\hat c^0}(\eta) \pm \xi\, P_{\hat c^1}(\eta) = T(\pm \xi)$.
Thus $w$ agrees with the codeword $\ev_{L_j}(T) \in \Cc_j$ on every fibre over $L'$, so $\Delta^{\mathrm{fib}}_j(w, \Cc_j) \le \delta$ by \cref{lem:fibre-dist}(2), a contradiction.
\end{proof}

Compared with \KB{Lemma~7.6}, the polynomial $T$ is built from the two decoded codewords directly: the generators of the line are the even and odd parts, so no change of generators is needed.

\begin{lemma}[Fibre errors survive folding]\label{lem:survive}
Let $j \in \iv{0}{\ell-1}$, $w, c \in \F^{L_j}$, and $\eta \in L_{j+1}$ such that $w$ and $c$ differ somewhere on the fibre over $\eta$. Then there is at most one $\theta \in \F$ with $\cfold_\theta(w)(\eta) = \cfold_\theta(c)(\eta)$.
\end{lemma}

\begin{proof}
Let $y = w - c$. By \cref{lem:cword}(1) and (2), $\cfold_\theta(w)(\eta) - \cfold_\theta(c)(\eta) = \cfold_\theta(y)(\eta) = y_e(\eta) + \theta\, y_o(\eta)$.
This is a polynomial of degree $\le 1$ in $\theta$. If it were identically zero, then $y_e(\eta) = y_o(\eta) = 0$, and $y$ would vanish on the fibre over $\eta$ (\cref{lem:cword}(2)). Hence it has at most one root.
\end{proof}

\subsection{Good challenges}\label{sec:sound:good}

\begin{definition}[Good challenge]\label{def:good}
Let $j \in \iv{1}{\ell}$, $w \in \F^{L_{j-1}}$ and $\theta \in \F$. We say that $\theta$ is \emph{good for $w$ at level $j$}, and write $\mathsf{Good}_j(w,\theta)$, if:
\begin{enumerate}
  \item in case $\Delta^{\mathrm{fib}}_{j-1}(w, \Cc_{j-1}) > \delta$: $\Delta(\cfold_\theta(w), \Cc_{j}) > \delta$;
  \item in case $\Delta^{\mathrm{fib}}_{j-1}(w, \Cc_{j-1}) \le \delta$, with $c = \mathrm{dec}_{j-1}(w)$: for every $\eta \in L_j$ such that $w$ and $c$ differ somewhere on the fibre over $\eta$, $\cfold_\theta(w)(\eta) \ne \cfold_\theta(c)(\eta)$.
\end{enumerate}
\end{definition}

\begin{lemma}[Few bad challenges]\label{lem:good}
For every $j \in \iv{1}{\ell}$ and $w \in \F^{L_{j-1}}$, at most $M_j$ values $\theta \in \F$ are not good for $w$ at level $j$.
\end{lemma}

\begin{proof}
In case 1, this is \cref{lem:far}. In case 2, by \cref{lem:survive}, each of the at most $M_j$ points $\eta$ over which $w$ and $c$ differ excludes at most one $\theta$.
\end{proof}

\begin{lemma}[One folding step]\label{lem:step}
Let $j \in \iv{1}{\ell}$, $w \in \F^{L_{j-1}}$, and $\theta \in \F$ good for $w$ at level $j$.
Let $c' \in \Cc_j$ and $Y = \{\eta \in L_j : \cfold_\theta(w)(\eta) = c'(\eta)\}$, and suppose $|Y| \ge (1-\delta) M_j$.
Then $\Delta^{\mathrm{fib}}_{j-1}(w, \Cc_{j-1}) \le \delta$; the codeword $c = \mathrm{dec}_{j-1}(w)$ satisfies $\cfold_\theta(c) = c'$; and $w = c$ on the fibre over every $\eta \in Y$.
\end{lemma}

\begin{proof}
We have $\Delta(\cfold_\theta(w), c') \le \delta$, so case 1 of \cref{def:good} is excluded: $\Delta^{\mathrm{fib}}_{j-1}(w, \Cc_{j-1}) \le \delta$.
By \cref{lem:cword}(4) (with $2N_j = N_{j-1} \le M_{j-1}$), $\cfold_\theta(c) \in \Cc_j$.
By locality (\cref{lem:cword}(1)), $\cfold_\theta(c)$ agrees with $\cfold_\theta(w)$ at every $\eta$ over whose fibre $w = c$, hence on at least $(1-\delta)M_j$ points (\cref{lem:fibre-dist}(2)); and $\cfold_\theta(w)$ agrees with $c'$ on $Y$.
By \cref{lem:metric}, $\cfold_\theta(c)$ and $c'$ agree on at least $(1-2\delta)M_j \ge \rho M_j = N_j$ points, so they are equal (\cref{lem:rs-params}(2)).
Finally, for $\eta \in Y$, $\cfold_\theta(w)(\eta) = c'(\eta) = \cfold_\theta(c)(\eta)$, so by case 2 of \cref{def:good}, $w = c$ on the fibre over $\eta$.
\end{proof}

\subsection{The soundness theorem}\label{sec:sound:thm}

Fix a deterministic prover $\mathsf{P}^*$ (\cref{lem:randomised} extends the results to randomised provers) and an instance $(w_0; z, v)$ with $w_0 \in \F^L$.
The prover's messages are functions of the earlier challenges: $s_j$ of $\theta_{\le j-1}$, the oracle $w_j$ ($j \in \iv{1}{\ell-1}$) of $\theta_{\le j}$, and $g$ of $\theta_{\le \ell}$.
As in \cref{sec:scheme:protocol}, $G = V_g$ and $w_\ell = \ev_{L_\ell}(G) \in \Cc_\ell$.

\paragraph{Passing sets.}
For fixed challenges $\theta \in \F^\ell$, let $\Pp_\ell = L_\ell$ and, for $j = \ell-1, \dots, 0$,
\[
  \Pp_j \;=\; \bigl\{\, \xi \in L_j \;:\; \xi^2 \in \Pp_{j+1} \text{ and } \cfold_{\theta_{j+1}}(w_j)(\xi^2) = w_{j+1}(\xi^2) \,\bigr\}.
\]
Let $\pi = \pi(\theta) = |\Pp_0|/M$.

\begin{lemma}[Passing sets, \KB{Lemma~7.11}]\label{lem:passing}
For fixed $\theta \in \F^\ell$:
\begin{enumerate}
  \item for $j \in \iv{0}{\ell}$, a point $\xi_0 \in L$ has $\xi_0^{2^j} \in \Pp_j$ if and only if the query checks at levels $j+1,\dots,\ell$ pass for $\xi_0$; in particular $\xi_0$ passes all the query checks if and only if $\xi_0 \in \Pp_0$;
  \item for $j < \ell$, $\Pp_j$ is a union of fibres, and the squaring map sends it two-to-one onto a set $\Pp_j^2 \subseteq \Pp_{j+1}$;
  \item $|\Pp_j| \ge \pi M_j$ for every $j \in \iv{0}{\ell}$.
\end{enumerate}
\end{lemma}

\begin{proof}
(1) Downward induction on $j$, from the recursive definition. (2) The defining condition of $\Pp_j$ depends only on $\xi^2$, and squaring is two-to-one on $L_j$ (\cref{lem:power}). (3) By (2), $|\Pp_j|/M_j = |\Pp_j^2|/M_{j+1} \le |\Pp_{j+1}|/M_{j+1}$; iterate from $j = 0$.
\end{proof}

\begin{lemma}[The codeword chain]\label{lem:chain}
Fix $\theta \in \F^\ell$ such that $\theta_j$ is good for $w_{j-1}$ at level $j$ for every $j \in \iv{1}{\ell}$, and $\pi > 1-\delta$.
Then $\Delta^{\mathrm{fib}}_j(w_j, \Cc_j) \le \delta$ for every $j \in \iv{0}{\ell-1}$, and the codewords $c_j = \mathrm{dec}_j(w_j)$ for $j < \ell$ and $c_\ell = w_\ell$ satisfy $\cfold_{\theta_{j+1}}(c_j) = c_{j+1}$ and $w_j = c_j$ on $\Pp_j$, for every $j \in \iv{0}{\ell-1}$.
Moreover, $g = \mathrm{cres}_{\theta}(f^*)$, where $f^*$ is the decoded table of $w_0$.
\end{lemma}

\begin{proof}
We show by downward induction on $j \in \iv{0}{\ell}$ that $w_j = c_j$ on $\Pp_j$, together with the other claims for $j < \ell$. For $j = \ell$ this is trivial.
Let $j < \ell$. For $\eta \in \Pp_j^2 \subseteq \Pp_{j+1}$, $\cfold_{\theta_{j+1}}(w_j)(\eta) = w_{j+1}(\eta) = c_{j+1}(\eta)$ by the definition of $\Pp_j$ and the induction hypothesis.
So the set $Y$ of \cref{lem:step} (with $w = w_j$, $\theta = \theta_{j+1}$, $c' = c_{j+1}$, at level $j+1$) contains $\Pp_j^2$, of size $|\Pp_j|/2 \ge \pi M_{j+1} > (1-\delta)M_{j+1}$ (\cref{lem:passing}).
\Cref{lem:step} gives $\Delta^{\mathrm{fib}}_j(w_j,\Cc_j) \le \delta$, $\cfold_{\theta_{j+1}}(c_j) = c_{j+1}$, and $w_j = c_j$ on every fibre over $\Pp_j^2$, that is, on $\Pp_j$.

For the last claim, extend $\theta$ arbitrarily to a point of $\F^n$. Then $c_0 = \Enc(f^*) = \ev_L(V_{f^*})$, and by \cref{cor:cword-chain}, $c_\ell$, obtained from $c_0$ by $\ell$ classical folds, is $\ev_{L_\ell}(V_{\mathrm{cres}_\theta(f^*)})$.
Since $c_\ell = w_\ell = \ev_{L_\ell}(V_g)$ and both $V_g$ and $V_{\mathrm{cres}_\theta(f^*)}$ lie in $\F[X]_{<N_\ell}$ with $N_\ell \le M_\ell$, \cref{lem:rs-params}(1) gives $V_g = V_{\mathrm{cres}_\theta(f^*)}$, hence $g = \mathrm{cres}_\theta(f^*)$ by comparing coefficients.
\end{proof}

\begin{theorem}[Soundness]\label{thm:sound}
Let $1 \le \ell \le n$, and put
\[
  \varepsilon_{\mathrm{fold}} \;=\; \sum_{j=1}^{\ell} \frac{M_j}{|\F|} \;=\; \frac{M\,(1 - 2^{-\ell})}{|\F|} \;<\; \frac{M}{|\F|}.
\]
\begin{enumerate}
  \item If $\Delta^{\mathrm{fib}}_0(w_0, \Cc_0) > \delta$, then $\mathsf{V}$ accepts with probability at most $\varepsilon_{\mathrm{fold}} + (1-\delta)^{\kappa}$.
  \item If $\Delta^{\mathrm{fib}}_0(w_0, \Cc_0) \le \delta$ and the decoded table $f^*$ of $w_0$ satisfies $P_{f^*}(z) \ne v$, then $\mathsf{V}$ accepts with probability at most
  \[
    \varepsilon_{\mathrm{fold}} \;+\; \frac{\ell}{|\F|} \;+\; (1-\delta)^{\kappa}.
  \]
\end{enumerate}
\end{theorem}

\begin{proof}
\emph{Honest values.} In case 2, let $f^*_j = \mathrm{cres}_{\theta_{\le j}}(f^*)$ be the tables of the honest prover for $f^*$, let $s^*_j = h_{f^*_{j-1}, z_{>j}}$ be its round polynomials \eqref{eq:honest-sj}, and let $v^*_j = P_{f^*_j}(z_{>j})$ for $j \in \iv{0}{\ell}$.
By \cref{lem:honest}, $v^*_0 = P_{f^*}(z)$, $s^*_j(z_j) = v^*_{j-1}$ and $s^*_j(\theta_j) = v^*_j$ for $j \in \iv{1}{\ell}$, and $v^*_\ell = P_{f^*_\ell}(z_{>\ell})$.
Note that $s^*_j$ depends on $\theta_{\le j-1}$ only.

\emph{Bad events.} For $j \in \iv{1}{\ell}$, let $E_j$ be the set of $\theta_{\le j}$ such that $\theta_j$ is not good for $w_{j-1}$ at level $j$, or (in case 2) $s_j \ne s^*_j$ and $s_j(\theta_j) = s^*_j(\theta_j)$.
For fixed $\theta_{\le j-1}$, the word $w_{j-1}$ and the polynomials $s_j, s^*_j \in \F[T]_{<2}$ are fixed, so at most $M_j$ values of $\theta_j$ (\cref{lem:good}) plus, in case 2, at most one value (\cref{lem:round-poly}(3)) complete $\theta_{\le j-1}$ into $E_j$.
By \cref{lem:sequential}, $\Pr[\exists j : \theta_{\le j} \in E_j] \le \varepsilon_{\mathrm{fold}}$ in case 1 and $\le \varepsilon_{\mathrm{fold}} + \ell/|\F|$ in case 2.

\emph{No bad event.} Suppose that no $E_j$ occurs and that $\mathsf{V}$ accepts. We show that $\pi \le 1-\delta$.
Assume instead $\pi > 1-\delta$. Then \cref{lem:chain} applies. In case 1 it gives $\Delta^{\mathrm{fib}}_0(w_0,\Cc_0) \le \delta$, a contradiction.
In case 2, we show $v_j \ne v^*_j$ by induction on $j \in \iv{0}{\ell}$.
For $j = 0$, $v_0 = v \ne P_{f^*}(z) = v^*_0$.
If $v_{j-1} \ne v^*_{j-1}$, the restriction check of index $j$ gives $s_j(z_j) = v_{j-1} \ne v^*_{j-1} = s^*_j(z_j)$, so $s_j \ne s^*_j$; since $\theta_{\le j} \notin E_j$, $v_j = s_j(\theta_j) \ne s^*_j(\theta_j) = v^*_j$.
But \cref{lem:chain} gives $g = f^*_\ell$, so the closure check reads $v_\ell = P_g(z_{>\ell}) = v^*_\ell$, a contradiction.

\emph{Conclusion.} Acceptance implies that some $E_j$ occurs, or that $\pi(\theta) \le 1-\delta$ and all $\kappa$ query points lie in $\Pp_0$ (\cref{lem:passing}(1)).
By \cref{lem:queries} with $\Omega' = \F^\ell$, $\mathcal{A} = \{\theta : \pi(\theta) \le 1-\delta\}$ and $\mathcal{S}(\theta) = \Pp_0$, the second event has probability at most $(1-\delta)^\kappa$.
\end{proof}

The bound of case 2 holds without the hypothesis $\Delta^{\mathrm{fib}}_0(w_0, \Cc_0) \le \delta$, with $f^*$ the zero table when there is no decoded table; the Lean statement is in this form.
Compared with \KB{Theorem~7.13}, the term $2\ell/|\F|$ of the degree-$2$ sumcheck becomes $\ell/|\F|$.

\begin{corollary}[Evaluation binding]\label{cor:binding}
Let $1 \le \ell \le n$ and $\varepsilon = \varepsilon_{\mathrm{fold}} + \ell/|\F| + (1-\delta)^{\kappa}$.
For every $w_0 \in \F^L$ and $z \in \F^n$, if some prover makes $\mathsf{V}$ accept the instance $(w_0; z,v)$ with probability greater than $\varepsilon$, then $\Delta^{\mathrm{fib}}_0(w_0,\Cc_0) \le \delta$ and $v = P_{f^*}(z)$ for the decoded table $f^*$ of $w_0$.
In particular, $\Pieval$ is $\varepsilon$-evaluation binding (\cref{def:binding}).
\end{corollary}

\begin{proof}
Immediate from \cref{thm:sound} and \cref{lem:randomised}; $f^*$ depends on $w_0$ only.
\end{proof}

\begin{corollary}[The sum protocol]\label{cor:sum}
Let $1 \le \ell \le n$ and $\varepsilon$ as in \cref{cor:binding}.
If $\Delta^{\mathrm{fib}}_0(w_0, \Cc_0) \le \delta$ and the decoded table $f^*$ of $w_0$ satisfies $\sum_a f^*(a) \ne v$, the verifier of $\Pisum$ accepts $(w_0; v)$ with probability at most $\varepsilon$.
If some prover makes it accept with probability greater than $\varepsilon$, then $\Delta^{\mathrm{fib}}_0(w_0,\Cc_0) \le \delta$ and $v = \sum_a f^*(a)$.
\end{corollary}

\begin{proof}
$\Pisum$ is $\Pieval$ at $z = \mathbf 1$, and $P_{f^*}(\mathbf 1) = \sum_a f^*(a)$ by \eqref{eq:sum-is-eval}; apply \cref{thm:sound}(2) and \cref{cor:binding}.
\end{proof}

\subsection{The case \texorpdfstring{$\ell = 0$}{l = 0}}\label{sec:sound:l0}

For $\ell = 0$ there is no fold, and the fibre distance is replaced by the Hamming distance.

\begin{proposition}[No folding]\label{prop:l0}
Let $\ell = 0$. If a deterministic prover makes $\mathsf{V}$ accept $(w_0; z,v)$ with probability greater than $(1-\delta)^{\kappa}$, then $\Delta(w_0, \Cc_0) < \delta$, and $v = P_{f^H}(z)$, where $f^H$ is the table such that $\Enc(f^H)$ is the unique codeword within Hamming distance $\delta$ of $w_0$ (\cref{lem:rs-params}(3)).
\end{proposition}

\begin{proof}
For $\ell = 0$ the prover sends $g$ before any challenge, so $g$ and $G = V_g$ are fixed, and the verifier accepts if and only if $v = P_g(z)$ and all query points lie in $S = \{\xi \in L : w_0(\xi) = G(\xi)\}$.
This has probability $0$ or $(|S|/M)^\kappa$. If it exceeds $(1-\delta)^\kappa$, then $v = P_g(z)$ and $|S| > (1-\delta)M$, so $\Delta(w_0, \ev_L(G)) < \delta$; hence $\ev_L(G) = \Enc(g)$ is the unique codeword within distance $\delta$ of $w_0$, and $g = f^H$ (\cref{lem:encoding}).
\end{proof}

\subsection{Round-by-round knowledge soundness}\label{sec:sound:rbr}

We now prove round-by-round knowledge soundness (\cref{def:rbr-knowledge}); \cref{sec:sound:pq} derives from it the relaxed notion of \cref{def:rrbr}, which is the hypothesis of \cref{thm:cdhz}.

\paragraph{Relation.}
Let $1 \le \ell \le n$.
We view $\Pieval$ as an IOP for the relation $\mathcal{R}^{\delta}_{\mathrm{eval}}$ of \cref{sec:proofs:pcs}, with $\mathrm{dist}(w, \Enc(f)) = \Delta^{\mathrm{fib}}_0(w, \Enc(f))$; condition~\eqref{eq:dist-unique} holds by \cref{lem:unique} and the injectivity of $\Enc$.
A table $f$ is a witness for $(w_0;z,v)$ if and only if $\Delta^{\mathrm{fib}}_0(w_0, \Cc_0) \le \delta$, $f$ is the decoded table of $w_0$, and $P_f(z) = v$.

\paragraph{Notation.}
The rounds are those of \cref{sec:scheme:protocol}, and $\tau_j$ is the partial transcript after $j$ rounds.
For $j \in \iv{1}{\ell}$ let $\hat w_j = \cfold_{\theta_j}(w_{j-1}) \in \F^{L_j}$, which is determined by $\tau_j$, and for $j \in \iv{1}{\ell-1}$ let $\mathcal{B}_j = \{\eta \in L_j : \hat w_j(\eta) \ne w_j(\eta)\}$, which is determined by $\tau_{j+1}$.
The \emph{restriction check of index $j$} is the equation $s_j(z_j) = v_{j-1}$, which is determined by $\tau_j$.
For $c \in \Cc_j$, let $\lambda[c]$ be the coefficient vector, with $n-j$ variables, of $P_c$: $\lambda[c](a) = \coef{a}{P_c}$ for $a \in \iv{0}{N_j - 1}$, so that $c = \ev_{L_j}(V_{\lambda[c]})$.

\paragraph{Witness sets and claimed values.}
For $c \in \Cc_0$ let $\mathcal{W}_0(c) = \{\xi \in L : w_0(\xi) = c(\xi) \text{ and } w_0(-\xi) = c(-\xi)\}$, and for $j \in \iv{1}{\ell}$ and $c \in \Cc_j$ let
\[
  \mathcal{W}_j(c) \;=\; \bigl\{\, \xi \in L \;:\; \xi^{2^i} \notin \mathcal{B}_i \text{ for all } i \in \iv{1}{j-1}, \text{ and } \hat w_j(\xi^{2^j}) = c(\xi^{2^j}) \,\bigr\},
\]
which is determined by $\tau_j$. Let $\mathrm{dens}_j(c) = |\mathcal{W}_j(c)|/M$, and for $j \in \iv{0}{\ell}$ and $c \in \Cc_j$ let
\[
  \sigma_j(c) \;=\; P_{\lambda[c]}(z_{j+1}, \dots, z_{n}),
\]
the value that the running claim $v_j$ takes if the prover is honest with respect to $c$ (\cref{lem:honest}).
Unlike in KBFold (\KB{\S7.6}), $\sigma_j(c)$ carries no factor depending on the challenges.

\begin{definition}[Doomed set and extractor]\label{def:state}
\begin{itemize}
  \item The empty transcript $\tau_0$ is doomed.
  \item For $j \in \iv{1}{\ell}$, the partial transcript $\tau_j$ is \emph{not doomed} if the restriction checks of indices $1,\dots,j$ hold and there is $c \in \Cc_j$ with
  \[
    \mathrm{dens}_j(c) \;\ge\; 1 - \delta
    \qquad\text{and}\qquad
    v_j \;=\; \sigma_j(c);
  \]
  otherwise it is \emph{doomed}.
  \item A full transcript $\tau_{\ell+1}$ is doomed if and only if the verifier rejects it.
\end{itemize}
Let $\Dd$ be the set of doomed transcripts.
The \emph{extractor} $\mathsf{Ext}$, on every input $(\tau, \mathsf{m})$, reads the oracle part $w_0$ of the instance of $\tau$, and outputs the decoded table of $w_0$ if $\Delta^{\mathrm{fib}}_0(w_0,\Cc_0) \le \delta$, and the zero table otherwise.
\end{definition}

By \cref{rem:ext}, $\mathsf{Ext}$ runs in time polynomial in $M$.
Its output is a witness for $(w_0;z,v)$ if and only if $\Delta^{\mathrm{fib}}_0(w_0,\Cc_0) \le \delta$ and the decoded table $f^*$ satisfies $P_{f^*}(z) = v$: if $\Delta^{\mathrm{fib}}_0(w_0,\Cc_0) > \delta$, no table is a witness.

\begin{lemma}[One round of the chain, \KB{Lemma~7.17}]\label{lem:rbr-step}
Let $j \in \iv{1}{\ell}$, fix $\tau_{j-1}$ and the round-$j$ message (hence $w_{j-1}$), and let $\theta_j$ be good for $w_{j-1}$ at level $j$.
If $c' \in \Cc_j$ satisfies $\mathrm{dens}_{j}(c') \ge 1-\delta$, then $\Delta^{\mathrm{fib}}_{j-1}(w_{j-1}, \Cc_{j-1}) \le \delta$, the codeword $c = \mathrm{dec}_{j-1}(w_{j-1})$ satisfies $\cfold_{\theta_j}(c) = c'$, and $\mathcal{W}_{j}(c') \subseteq \mathcal{W}_{j-1}(c)$; in particular $\mathrm{dens}_{j-1}(c) \ge 1-\delta$.
\end{lemma}

\begin{proof}
Let $Y = \{\eta \in L_{j} : \hat w_j(\eta) = c'(\eta)\}$. Every $\xi \in \mathcal{W}_{j}(c')$ has $\xi^{2^j} \in Y$, and the map $\xi \mapsto \xi^{2^j}$ is $2^j$-to-$1$ from $L$ onto $L_j$ (\cref{lem:power}); so $|Y| \ge |\mathcal{W}_j(c')|/2^j \ge (1-\delta)M_j$.
\Cref{lem:step} gives the first two claims, and $w_{j-1} = c$ on the fibre over every $\eta \in Y$.
Let $\xi \in \mathcal{W}_{j}(c')$, and $\eta = \xi^{2^j} \in Y$; the fibre over $\eta$ is $\{\pm \xi^{2^{j-1}}\}$, on which $w_{j-1} = c$.
If $j = 1$, this says $\xi \in \mathcal{W}_0(c)$.
If $j \ge 2$, the condition $\xi^{2^{j-1}} \notin \mathcal{B}_{j-1}$ (part of the definition of $\mathcal{W}_{j}(c')$) gives $\hat w_{j-1}(\xi^{2^{j-1}}) = w_{j-1}(\xi^{2^{j-1}}) = c(\xi^{2^{j-1}})$, and the conditions $\xi^{2^i} \notin \mathcal{B}_i$ for $i \le j-2$ are common to both sets; hence $\xi \in \mathcal{W}_{j-1}(c)$.
\end{proof}

\begin{lemma}[Restriction step for a codeword]\label{lem:rbr-restr}
Let $j \in \iv{1}{\ell}$ and $c \in \Cc_{j-1}$, and let $s^{c}_j = h_{\lambda[c],\, z_{>j}} \in \F[T]_{<2}$ (\cref{def:round-poly}).
Then $s^{c}_j(z_j) = \sigma_{j-1}(c)$, and $s^{c}_j(\theta) = \sigma_{j}(\cfold_\theta(c))$ for every $\theta \in \F$.
\end{lemma}

\begin{proof}
By \cref{lem:round-poly}(1), $s^c_j(t) = P_{\lambda[c]}(t, z_{>j})$ for every $t \in \F$; for $t = z_j$ this is $\sigma_{j-1}(c)$.
Since $c = \ev_{L_{j-1}}(V_{\lambda[c]})$ and $V_{\lambda[c]} \in \F[X]_{<N_{j-1}}$ with $N_{j-1} = 2N_j \le M_{j-1}$, \cref{lem:cword}(4) and \cref{thm:dictionary} give $\cfold_\theta(c) = \ev_{L_j}(\cfold_\theta(V_{\lambda[c]})) = \ev_{L_j}(V_{\mathrm{cres}_\theta(\lambda[c])})$, so $\lambda[\cfold_\theta(c)] = \mathrm{cres}_\theta(\lambda[c])$ by \cref{lem:rs-params}(1).
By \cref{lem:cres}(1), $\sigma_j(\cfold_\theta(c)) = P_{\mathrm{cres}_\theta(\lambda[c])}(z_{>j}) = P_{\lambda[c]}(\theta, z_{>j}) = s^c_j(\theta)$.
\end{proof}

\begin{theorem}[Round-by-round knowledge soundness]\label{thm:rbr}
Let $1 \le \ell \le n$, and let $\Dd$ and $\mathsf{Ext}$ be as in \cref{def:state}.
\begin{enumerate}
  \item (Round $1$.) If the output of $\mathsf{Ext}$ is not a witness for $(w_0; z,v)$, then for every round-$1$ message $s_1$, $\tau_1$ is not doomed with probability at most $(M_1 + 1)/|\F|$ over $\theta_1$.
  \item (Round $j$, $2 \le j \le \ell$.) If $\tau_{j-1}$ is doomed, then for every round-$j$ message $(s_j, w_{j-1})$, $\tau_j$ is not doomed with probability at most $(M_j+1)/|\F|$ over $\theta_j$.
  \item (Round $\ell+1$.) If $\tau_\ell$ is doomed, then for every final message $g$, the verifier accepts with probability at most $(1-\delta)^{\kappa}$ over the query points.
\end{enumerate}
Consequently, $\Pieval$ is round-by-round knowledge sound for $\mathcal{R}^{\delta}_{\mathrm{eval}}$ (\cref{def:rbr-knowledge}) with errors $\varepsilon_j = (M_j+1)/|\F|$ for $j \in \iv{1}{\ell}$ and $\varepsilon_{\ell+1} = (1-\delta)^\kappa$, and in particular with
\[
  \varepsilon_{\mathrm{rbr}} \;=\; \max\Bigl\{ \frac{M_1 + 1}{|\F|},\; (1-\delta)^{\kappa} \Bigr\}.
\]
\end{theorem}

\begin{proof}
\emph{Items 1 and 2.} Fix $j \in \iv{1}{\ell}$, the instance, $\tau_{j-1}$ and the round-$j$ message.
Then $w_{j-1}$, $s_j$ and, if $\Delta^{\mathrm{fib}}_{j-1}(w_{j-1},\Cc_{j-1}) \le \delta$, the codeword $c = \mathrm{dec}_{j-1}(w_{j-1})$ and the polynomial $s^{c}_j$ of \cref{lem:rbr-restr} are fixed.
Call $\theta_j$ \emph{bad} if it is not good for $w_{j-1}$ at level $j$, or if $\Delta^{\mathrm{fib}}_{j-1}(w_{j-1},\Cc_{j-1}) \le \delta$, $s_j \ne s^{c}_j$ and $s_j(\theta_j) = s^{c}_j(\theta_j)$.
By \cref{lem:good} and \cref{lem:round-poly}(3), at most $M_j + 1$ values of $\theta_j$ are bad.
We show that if $\theta_j$ is not bad and $\tau_j$ is not doomed, then the hypothesis of the item fails.

Let $c' \in \Cc_j$ witness that $\tau_j$ is not doomed: $\mathrm{dens}_j(c') \ge 1-\delta$, $v_j = \sigma_j(c')$, and the restriction checks of indices $1,\dots,j$ hold.
By \cref{lem:rbr-step}, $c$ exists, $\cfold_{\theta_j}(c) = c'$ and $\mathrm{dens}_{j-1}(c) \ge 1-\delta$.
By \cref{lem:rbr-restr}, $s_j(\theta_j) = v_j = \sigma_j(\cfold_{\theta_j}(c)) = s^{c}_j(\theta_j)$, so $s_j = s^{c}_j$ because $\theta_j$ is not bad.
The restriction check of index $j$ then gives $v_{j-1} = s_j(z_j) = s^{c}_j(z_j) = \sigma_{j-1}(c)$.

For $j \ge 2$, the codeword $c$ therefore witnesses that $\tau_{j-1}$ is not doomed, contradicting the hypothesis of item 2.
For $j = 1$, $\Delta^{\mathrm{fib}}_0(w_0,\Cc_0) \le \delta$ and $c = \mathrm{dec}_0(w_0) = \Enc(f^*)$ for the decoded table $f^*$, so $\lambda[c] = f^*$ and $v = v_0 = \sigma_0(c) = P_{f^*}(z)$; thus the output $f^*$ of $\mathsf{Ext}$ is a witness, contradicting the hypothesis of item 1.

\emph{Item 3.} If a restriction check or the closure check fails, the verifier rejects. Otherwise, since $\lambda[w_\ell] = g$, the closure check reads $v_\ell = \sigma_\ell(w_\ell)$; as $\tau_\ell$ is doomed, $\mathrm{dens}_\ell(w_\ell) < 1-\delta$.
A query point $\xi$ passes all its checks exactly when $\hat w_i(\xi^{2^i}) = w_i(\xi^{2^i})$ for all $i \in \iv{1}{\ell}$, that is, when $\xi \in \mathcal{W}_\ell(w_\ell)$.
The $\kappa$ query points are independent and uniform, so the verifier accepts with probability at most $\mathrm{dens}_\ell(w_\ell)^{\kappa} < (1-\delta)^{\kappa}$ (\cref{lem:queries}).

\emph{Round-by-round knowledge soundness.} We check the three conditions of \cref{def:rbr-knowledge}.
Condition 1 holds because $\tau_0$ is doomed, and condition 3 because a full transcript is doomed exactly when it is rejected.
For condition 2 in round $1$, item 1 says that an escape probability above $\varepsilon_1$ forces $\mathsf{Ext}$ to output a witness.
In rounds $2,\dots,\ell+1$, items 2 and 3 bound the escape probability from a doomed transcript by $\varepsilon_j$ outright.
Finally, $M_j \le M_1$ for all $j \ge 1$.
\end{proof}

\begin{remark}[Binding from knowledge soundness]\label{rem:bgtz}
\Cref{def:rbr-knowledge} is the definition of Block, Garreta, Tiwari and Zaj\k{a}c~\cite{BGTZ24}, with the query positions placed in the challenge of the last round (\KB{Remark~7.20}).
Because the relation $\mathcal{R}^{\delta}_{\mathrm{eval}}$ requires the witness to be the table decoded from the commitment, the extracted witness is bound to the commitment; by \cref{lem:rbr-binding}, \cref{thm:rbr} implies $\varepsilon$-evaluation binding with $\varepsilon = \sum_{j=1}^{\ell}(M_j+1)/|\F| + (1-\delta)^\kappa$, which is also the bound of \cref{cor:binding}.
For the sum protocol, the same holds for $\mathcal{R}^{\delta}_{\mathrm{sum}}$ by \eqref{eq:sum-is-eval}.
\end{remark}

\begin{proposition}[Round-by-round knowledge soundness for $\ell = 0$]\label{prop:rbr-l0}
Let $\ell = 0$, and define $\mathcal{R}^{\delta}_{\mathrm{eval}}$ with $\mathrm{dist} = \Delta$; condition~\eqref{eq:dist-unique} holds by \cref{lem:rs-params}(3) and the injectivity of $\Enc$ (\cref{lem:encoding}).
Let $\mathsf{Ext}$ output the table $f^H$ of \cref{prop:l0} if $\Delta(w_0, \Cc_0) \le \delta$ (computed with \cref{fact:decode} and an inverse NTT), and the zero table otherwise, and let $\Dd$ consist of $\tau_0$ and the rejected full transcripts.
Then $\Pieval$ is round-by-round knowledge sound for $\mathcal{R}^{\delta}_{\mathrm{eval}}$ with error $\varepsilon_1 = (1-\delta)^\kappa$.
\end{proposition}

\begin{proof}
Conditions 1 and 3 of \cref{def:rbr-knowledge} hold by construction. For condition 2, fix the message $g$: if the acceptance probability exceeds $(1-\delta)^\kappa$, the proof of \cref{prop:l0} gives $\Delta(w_0, \Cc_0) < \delta$ and $v = P_{f^H}(z)$, so the output of $\mathsf{Ext}$ is a witness.
\end{proof}

\subsection{The non-interactive scheme}\label{sec:sound:pq}

We now show that $\Pieval$, written as an interactive oracle reduction, satisfies \cref{def:rrbr}, and apply \cref{thm:cdhz}.
The encoding as an IOR is that of \KB{\S7.7}; the only changes are the shorter round messages and the restriction check.
Let $1 \le \ell \le n$ and $\delta^* = (1-\rho)/2$.
Everything proved so far in this section holds for every $\delta \in (0, \delta^*]$; below, $\delta$ varies in $[0,1]$.

\paragraph{Fibre strings.}
For each $j \in \iv{0}{\ell-1}$ and $\eta \in L_{j+1}$, fix one of the two square roots $\xi_\eta \in L_j$ of $\eta$ (\cref{lem:power}); the other is $-\xi_\eta$.
Let $\Sigma = \F^2$, and for $w \in \F^{L_j}$ let $\mathrm{fib}_j(w) \in \Sigma^{L_{j+1}}$ be the string $\eta \mapsto (w(\xi_\eta), w(-\xi_\eta))$, indexed by $L_{j+1}$ in a fixed order.
The map $\mathrm{fib}_j$ is a bijection from $\F^{L_j}$ onto $\Sigma^{M_{j+1}}$, and by \cref{def:fibre-dist},
\begin{equation}\label{eq:fib-hamming}
  \Delta\bigl(\mathrm{fib}_j(w), \mathrm{fib}_j(c)\bigr) \;=\; \Delta^{\mathrm{fib}}_j(w, c) \qquad (w, c \in \F^{L_j}).
\end{equation}
For a partial string $u$ over $\Sigma$, the \emph{filling} $\bar u$ replaces every $\bot$ by $(0,0)$, and $\mathcal{U}(u)$ is the set of positions where $u$ is $\bot$.
Since $\bar u$ and $u$ agree on $\mathrm{Dom}(u)$, $\Delta(\bar u, u') \le \Delta(u, u')$ for every full string $u'$.

\paragraph{The relation.}
Let
\[
  \mathcal{R}_{\mathrm{eval}} \;=\; \bigl\{\, ((z,v),\, y,\, f) \;:\; y = \mathrm{fib}_0(\Enc(f)) \ \text{ and } \ P_f(z) = v \,\bigr\},
\]
with explicit instance $(z, v) \in \F^n \times \F$, implicit instance $y \in \Sigma^{M/2}$ and witness a table $f$ over $\F$.
For a partial string $y$ and $\delta \in [0,1]$, $((z,v), y, f) \in (\mathcal{R}_{\mathrm{eval}})_{\le\delta}$ if and only if $\Delta(y, \mathrm{fib}_0(\Enc(f))) \le \delta$ and $P_f(z) = v$.
For a full string $y = \mathrm{fib}_0(w_0)$, by \eqref{eq:fib-hamming}, this says $((w_0;z,v),f) \in \mathcal{R}^{\delta}_{\mathrm{eval}}$ in the notation of \cref{sec:sound:rbr}.
The sum relation $\mathcal{R}_{\mathrm{sum}}$ is the restriction to $z = \mathbf 1$.

\paragraph{The protocol as an IOR.}
Let $\Pieval^{\Sigma}$ be $\Pieval$ with instance $((z,v); y)$, and with messages and challenges encoded as follows.
\begin{itemize}
  \item For $j \in \iv{1}{\ell}$, the round-$j$ message is the symbol $(a_j, b_j) \in \Sigma$ of the coefficients of $s_j(T) = a_j + b_j T$, followed, if $j \ge 2$, by the oracle part $\mathrm{fib}_{j-1}(w_{j-1})$. The final message is $g$, written as $\lceil N_\ell/2 \rceil$ symbols, padded with a zero if $N_\ell = 1$. Padding entries are ignored.
  \item Fix an integer $\beta \ge \log_2 |\F|$ and a bijection $\mathrm{num} : \iv{0}{|\F|-1} \to \F$. The challenge of round $j \in \iv{1}{\ell}$ is $\mathsf{vr}_j \in \{0,1\}^\beta$, and $\theta_j = \phi(\mathsf{vr}_j)$, where $\phi(\mathsf{vr}) = \mathrm{num}(\mathrm{int}(\mathsf{vr}) \bmod |\F|)$ and $\mathrm{int}$ reads a bit string as a binary integer.
  \item Fix a bijection $\mathrm{pos} : \{0,1\}^{n+R} \to L$ (recall $M = 2^{n+R}$). The challenge of round $\ell+1$ is $\mathsf{vr}_{\ell+1} \in \{0,1\}^{\kappa(n+R)}$, cut into $\kappa$ blocks of $n+R$ bits, whose images under $\mathrm{pos}$ are the query points.
  \item On partial strings, the verifier rejects if an entry that it reads is $\bot$. Otherwise it runs the verifier of $\Pieval$ on the words $w_0 = \mathrm{fib}_0^{-1}(\bar y)$ and $w_{j-1} = \mathrm{fib}_{j-1}^{-1}(\overline{\text{oracle part of round } j})$, the round polynomials and the final table read from the filled messages; it outputs $(\top, ())$ if that verifier accepts, and rejects otherwise.
\end{itemize}
Every string of the prescribed length over $\Sigma$ encodes a message of $\Pieval$, and filled strings are such strings.
The query point $\xi_0 = \xi^{(i)}_0$ reads, for each level $i' \in \iv{1}{\ell}$, the entry at position $\xi_0^{2^{i'}}$ of the string $y$ (if $i' = 1$) or of the oracle part of round $i'$ (if $i' \ge 2$): this entry contains the two values $w_{i'-1}(\pm\xi_0^{2^{i'-1}})$ of the fold check of level $i'$, and, for $i' \ge 2$, one of them is the value $w_{i'-1}(\xi_0^{2^{i'-1}})$ compared at level $i'-1$.
The positions read are determined by the challenges.
So $\Pieval^{\Sigma}$ is an IOR from $\mathcal{R}_{\mathrm{eval}}$ to $\mathcal{R}_{\mathrm{acc}}$ with $\nu = \ell + 1$ rounds, $\beta_{\min} = \min(\beta, \kappa(n+R))$ and $l_{\max} = M/2$.
The compressed messages of \cref{rem:compress} are a bijective re-encoding of the messages that pass the restriction check, and the analysis below applies to them unchanged.

\begin{lemma}[Sampling, \KB{Lemma~7.22}]\label{lem:sampling}
For every $S \subseteq \F$, $\Pr_{\mathsf{vr} \leftarrow \{0,1\}^\beta}[\phi(\mathsf{vr}) \in S] \le |S|\,\bigl(1/|\F| + 2^{-\beta}\bigr)$.
The query points of round $\ell+1$ are independent and uniform in $L$.
\end{lemma}

\begin{proof}
Each $a \in \F$ has at most $\lceil 2^\beta/|\F| \rceil < 2^\beta/|\F| + 1$ preimages under $\phi$; divide by $2^\beta$ and sum over $S$.
The map $\{0,1\}^{\kappa(n+R)} \to L^{\times\kappa}$ induced by $\mathrm{pos}$ is a bijection.
\end{proof}

\paragraph{Erasures.}
For a statement and messages given as partial strings, let $\mathcal{U}_0 = \mathcal{U}(y) \subseteq L_1$ and, for $i \in \iv{1}{\ell-1}$, let $\mathcal{U}_i \subseteq L_{i+1}$ be the set of erased positions of the oracle part of round $i+1$.
We apply the notation of \cref{sec:sound:rbr} to the filled words $w_0, \dots, w_{\ell-1}$, round polynomials and final table defined above, and, for $j \in \iv{0}{\ell}$ and $c \in \Cc_j$, we put
\[
  \mathcal{W}^{\bot}_j(c) \;=\; \bigl\{\, \xi \in \mathcal{W}_j(c) \;:\; \xi^{2^{i}} \notin \mathcal{U}_{i-1} \text{ for all } i \in \iv{1}{\max(j,1)} \,\bigr\}, \qquad \mathrm{dens}^{\bot}_j(c) = |\mathcal{W}^{\bot}_j(c)|/M.
\]
For $\delta \in (0,\delta^*]$ and the statement at hand, let $\Dd^{\bot}_\delta$ be the set of partial transcripts obtained from \cref{def:state} with $\mathrm{dens}^{\bot}_j$ in place of $\mathrm{dens}_j$; it depends on the erasures of $y$ and of the messages, not only on the filled words.
Without erasures, $\mathcal{W}^{\bot}_j = \mathcal{W}_j$ and $\Dd^{\bot}_\delta$ is the doomed set of \cref{def:state}.

\begin{definition}[Knowledge state function of $\Pieval^{\Sigma}$]\label{def:kstate}
The algorithm $\mathsf{E}$, on every input, reads $y$, lets $w_0 = \mathrm{fib}_0^{-1}(\bar y)$, and outputs the decoded table of $w_0$ at radius $\delta^*$ (\cref{def:decoded} with $\delta = \delta^*$) if $\Delta^{\mathrm{fib}}_0(w_0, \Cc_0) \le \delta^*$, and the zero table otherwise.
For $\delta \in [0,1]$, a statement $((z,v), y)$, a transcript $\mathrm{tr}$ and a table $f$, let $\mathsf{KState}_\delta((z,v), y, \mathrm{tr}, f)$ be:
\begin{enumerate}
  \item $1$ if and only if $((z,v), y, f) \in (\mathcal{R}_{\mathrm{eval}})_{\le\delta}$, if $\mathrm{tr} = \mathrm{tr}_\emptyset$;
  \item $1$ if and only if the verifier of $\Pieval^{\Sigma}$ accepts, if $\mathrm{tr}$ is a full transcript;
  \item for $\mathrm{tr} = \tau_j$ the partial transcript after $j \in \iv{1}{\ell}$ rounds: $1$ if and only if $\tau_j \notin \Dd^{\bot}_\delta$, when $\delta \in (0, \delta^*]$; and $1$ when $\delta \notin (0, \delta^*]$;
  \item $\mathsf{KState}_\delta((z,v), y, \mathrm{tr}', f)$, if $\mathrm{tr} = \mathrm{tr}' \| \Pi$ ends with a prover message.
\end{enumerate}
\end{definition}

$\mathsf{E}$ runs in time polynomial in $M$ by \cref{rem:ext}, and it does not depend on $\delta$.
Conditions 1 and 3 of \cref{def:rrbr} hold by items 1 and 2, and condition 2 by item 4.

\begin{lemma}[Round-by-round steps with erasures]\label{lem:rbr-erasures}
Let $\delta \in (0, \delta^*]$, fix a statement given by a partial string $y$, and let $j \in \iv{1}{\ell+1}$.
Fix the partial transcript $\tau_{j-1}$ and the round-$j$ message (partial strings).
\begin{enumerate}
  \item If $j = 1$ and $((z,v), y, \mathsf{E}(y)) \notin (\mathcal{R}_{\mathrm{eval}})_{\le\delta}$, then at most $M_1 + 1$ values of $\theta_1$ make $\tau_1 \notin \Dd^{\bot}_\delta$.
  \item If $2 \le j \le \ell$ and $\tau_{j-1} \in \Dd^{\bot}_\delta$, then at most $M_j + 1$ values of $\theta_j$ make $\tau_j \notin \Dd^{\bot}_\delta$.
  \item If $j = \ell+1$ and $\tau_\ell \in \Dd^{\bot}_\delta$, then the verifier of $\Pieval^{\Sigma}$ accepts with probability at most $(1-\delta)^{\kappa}$ over the query points.
\end{enumerate}
\end{lemma}

\begin{proof}
We follow the proof of \cref{thm:rbr}, applied to the filled words, which are words of $\Pieval$.
First, \cref{lem:rbr-step} holds with $\mathcal{W}^{\bot}$: if $\xi \in \mathcal{W}^{\bot}_j(c')$, then $\xi \in \mathcal{W}_j(c')$, so $\xi \in \mathcal{W}_{j-1}(c)$ by \cref{lem:rbr-step}; the erasure conditions defining $\mathcal{W}^{\bot}_{j-1}(c)$, for $i \le \max(j-1,1) \le j$, are among those defining $\mathcal{W}^{\bot}_{j}(c')$. So $\mathcal{W}^{\bot}_j(c') \subseteq \mathcal{W}^{\bot}_{j-1}(c)$, and the size bound on $Y$ in that proof still holds because $\mathcal{W}^{\bot}_j(c') \subseteq \mathcal{W}_j(c')$.
\Cref{lem:rbr-restr} does not involve witness sets.

(1) and (2). Call $\theta_j$ \emph{bad} as in the proof of \cref{thm:rbr}; at most $M_j + 1$ values are bad. As there, if $\theta_j$ is not bad and $\tau_j \notin \Dd^{\bot}_\delta$, then $c = \mathrm{dec}_{j-1}(w_{j-1})$ exists, $\mathrm{dens}^{\bot}_{j-1}(c) \ge 1-\delta$ and $v_{j-1} = \sigma_{j-1}(c)$, and the restriction checks of indices $1, \dots, j-1$ hold.
For $j \ge 2$ this says $\tau_{j-1} \notin \Dd^{\bot}_\delta$, contradicting the hypothesis of (2).
For $j = 1$: $\mathcal{W}^{\bot}_0(c)$ is the set of $\xi \in L$ with $\xi^2 \notin \mathcal{U}_0$ and $w_0 = c$ on $\{\pm\xi\}$, so its size is twice the number of fibres $\eta \in L_1$ with $y(\eta) = \mathrm{fib}_0(c)(\eta)$ (in particular $y(\eta) \ne \bot$). So $\mathrm{dens}^{\bot}_0(c) \ge 1-\delta$ gives $\Delta(y, \mathrm{fib}_0(c)) \le \delta$.
Hence $\Delta^{\mathrm{fib}}_0(w_0, c) = \Delta(\bar y, \mathrm{fib}_0(c)) \le \delta \le \delta^*$, so $\mathsf{E}(y) = \lambda[c]$ by \cref{lem:unique} at radius $\delta^*$ and \cref{lem:encoding}; and $v = v_0 = \sigma_0(c) = P_{\lambda[c]}(z)$. Thus $((z,v), y, \mathsf{E}(y)) \in (\mathcal{R}_{\mathrm{eval}})_{\le\delta}$, contradicting the hypothesis of (1).

(3). If the verifier accepts, no entry it reads is $\bot$, its round polynomials and final table are those of the filled messages, and the checks of $\Pieval$ pass on the filled words. As in item 3 of \cref{thm:rbr}, the closure check gives $v_\ell = \sigma_\ell(w_\ell)$, so $\mathrm{dens}^{\bot}_\ell(w_\ell) < 1-\delta$ because $\tau_\ell \in \Dd^{\bot}_\delta$.
A query point $\xi$ that passes its checks lies in $\mathcal{W}_\ell(w_\ell)$, and it reads the positions $\xi^{2^{i}}$ of the strings with erasure sets $\mathcal{U}_{i-1}$, $i \in \iv{1}{\ell}$, which are therefore not erased; so $\xi \in \mathcal{W}^{\bot}_\ell(w_\ell)$.
By \cref{lem:sampling} and \cref{lem:queries}, the verifier accepts with probability at most $\mathrm{dens}^{\bot}_\ell(w_\ell)^\kappa < (1-\delta)^\kappa$.
\end{proof}

\begin{theorem}[Relaxed round-by-round knowledge soundness]\label{thm:rrbr}
$\Pieval^{\Sigma}$ has relaxed round-by-round knowledge soundness errors (\cref{def:rrbr}), with the knowledge state function and the algorithm $\mathsf{E}$ of \cref{def:kstate},
\[
  \varepsilon^{\mathrm{rr}}_j(\delta) \;=\; (M_j + 1)\Bigl(\frac{1}{|\F|} + \frac{1}{2^\beta}\Bigr) \quad (j \in \iv{1}{\ell}), \qquad \varepsilon^{\mathrm{rr}}_{\ell+1}(\delta) \;=\; (1-\delta)^{\kappa},
\]
for $\delta \in (0, \delta^*]$, and $\varepsilon^{\mathrm{rr}}_j(\delta) = 1$ for $\delta \notin (0, \delta^*]$.
\end{theorem}

\begin{proof}
For $\delta \notin (0,\delta^*]$ there is nothing to prove. Fix $\delta \in (0,\delta^*]$, a statement, a round $j \in \iv{1}{\ell+1}$ and $\mathrm{tr} = \tau_{j-1} \| \Pi_j$, where $\tau_0 = \mathrm{tr}_\emptyset$.
By item 4 of \cref{def:kstate}, $\mathsf{KState}_\delta(\cdot, \mathrm{tr}, \cdot) = \mathsf{KState}_\delta(\cdot, \tau_{j-1}, \cdot)$, and by items 2 and 3, $\mathsf{KState}_\delta(\cdot, \mathrm{tr}\|\mathsf{vr}_j, f)$ does not depend on $f$.
So the event of \cref{def:rrbr} is contained in the conjunction of $\mathsf{KState}_\delta(\cdot, \tau_{j-1}, \mathsf{E}(\cdot)) = 0$, which does not depend on $\mathsf{vr}_j$, and $\mathsf{KState}_\delta(\cdot, \tau_j, \cdot) = 1$.
If the first condition fails, the event is empty. Otherwise:
for $j = 1$, $((z,v), y, \mathsf{E}(y)) \notin (\mathcal{R}_{\mathrm{eval}})_{\le\delta}$; for $2 \le j \le \ell+1$, $\tau_{j-1} \in \Dd^{\bot}_\delta$.
For $j \le \ell$, \cref{lem:rbr-erasures}(1,2) and \cref{lem:sampling} with $|S| \le M_j + 1$ bound the probability that $\tau_j \notin \Dd^{\bot}_\delta$ by $\varepsilon^{\mathrm{rr}}_j(\delta)$.
For $j = \ell+1$, $\mathsf{KState}_\delta(\cdot, \tau_{\ell+1}, \cdot) = 1$ means acceptance, whose probability is at most $(1-\delta)^\kappa$ by \cref{lem:rbr-erasures}(3).
\end{proof}

\begin{corollary}[Post-quantum knowledge soundness]\label{cor:pq}
Let $\delta \in (0,\delta^*]$ and $\varepsilon_{\mathrm{rr}}(\delta) = \max\{(M_1+1)(1/|\F| + 2^{-\beta}),\ (1-\delta)^{\kappa}\}$.
The non-interactive scheme $\mathsf{BCS}[\Pieval^{\Sigma}]$ (\cref{sec:proofs:bcs}) satisfies the conclusion of \cref{thm:cdhz} with $\nu = \ell + 1$, $\beta_{\min} = \min(\beta, \kappa(n+R))$, $l_{\max} = M/2$ and this $\varepsilon_{\mathrm{rr}}(\delta)$.
In particular, for every quantum adversary that makes at most $Q$ queries to the random oracle and outputs a Merkle root $\mathrm{cm}$, a claim $(z, v)$ and a proof, except with probability $\varepsilon_{\mathrm{ext}}(Q)$, if the verifier accepts then the extractor outputs a partial string $y \in (\Sigma \sqcup \{\bot\})^{M/2}$, a table $f$ and a trapdoor such that:
the openings of all positions of $\mathrm{Dom}(y)$ are accepted under $\mathrm{cm}$; $y$ agrees with $\mathrm{fib}_0(\Enc(f))$ on at least $(1-\delta)M/2$ fibres; and $P_f(z) = v$.
For $\Pisum$ the last condition reads $\sum_a f(a) = v$.
\end{corollary}

\begin{proof}
\Cref{thm:rrbr} and \cref{thm:cdhz}; the last sentences unfold the relaxed committed relation $\widehat{\mathsf{CR}}[\mathcal{R}_{\mathrm{eval}},\delta]$ (\cref{sec:proofs:bcs}), $\bot$ entries counting as disagreements, and use \eqref{eq:sum-is-eval}.
\end{proof}

The extracted table is bound to the root: two extracted witnesses for the same root with different tables expose a hash collision.

\begin{lemma}[Uniqueness of relaxed committed witnesses, \KB{Lemma~7.27}]\label{lem:cr-unique}
Let $\delta \in (0,\delta^*]$ and let $\mathrm{cm}$ be a Merkle root. Let $\mathtt{cx} = ((z,v), \mathrm{cm})$ and $\mathtt{cx}' = ((z',v'), \mathrm{cm})$, and suppose that $(\mathtt{cx}, (y, f, \mathrm{td}))$ and $(\mathtt{cx}', (y', f', \mathrm{td}'))$ belong to $\widehat{\mathsf{CR}}[\mathcal{R}_{\mathrm{eval}}, \delta]$.
If $f \ne f'$, then $\mathrm{td}$ and $\mathrm{td}'$ yield two accepting openings, under $\mathrm{cm}$, of one position to different values, hence a collision of the random oracle.
In particular, if $z = z'$ and $v \ne v'$, such a collision is obtained.
\end{lemma}

\begin{proof}
Let $A = \{\eta \in L_1 : y(\eta) = \mathrm{fib}_0(\Enc(f))(\eta)\}$ and $A' = \{\eta \in L_1 : y'(\eta) = \mathrm{fib}_0(\Enc(f'))(\eta)\}$.
Then $A \subseteq \mathrm{Dom}(y)$ and $A' \subseteq \mathrm{Dom}(y')$, and $|A|, |A'| \ge (1-\delta)M/2$, so $|A \cap A'| \ge (1-2\delta)M/2 \ge \rho M/2 = N/2$.
If $y = y'$ on $A \cap A'$, then $\Enc(f)$ and $\Enc(f')$ agree on the at least $N$ points of the fibres over $A \cap A'$, hence are equal (\cref{lem:rs-params}(2)), and $f = f'$ (\cref{lem:encoding}).
So if $f \ne f'$, some $\eta \in A \cap A'$ has $y(\eta) \ne y'(\eta)$, and the openings of $\eta$ computed from $\mathrm{td}$ and $\mathrm{td}'$ are both accepted under $\mathrm{cm}$, with the same random oracle since the commitment index is fixed by the scheme; by the salted form of \KB{Lemma~3.13(1)} (\KB{\S3.5}), this yields a collision.
If $z = z'$ and $v \ne v'$, then $P_f(z) = v \ne v' = P_{f'}(z)$, so $f \ne f'$.
\end{proof}

\Cref{lem:cr-unique} is deterministic; for the interactive oracle proof, evaluation binding is \cref{cor:binding}.

\begin{remark}[Post-quantum parameters]\label{rem:pq-params}
Take the parameters of \KB{Remark~7.28}: $\rho = 1/4$, $\delta = \delta^* = 3/8$, $M \le 2^{32}$ (so $\ell \le n \le 30$), the quartic extension $\F = \mathbb{F}_p[\iota]/(\iota^4 - 7)$ of the Goldilocks field, of size $p^4 > 2^{255.99}$, $\beta = 320$, hash output length $\lambda_{\mathsf{H}} = 256$, and $\kappa = 248$ queries.
Then $(5/8)^{\kappa} < 2^{-168}$ and $(M_1+1)(1/|\F| + 2^{-\beta}) < 2^{-224}$, so $\varepsilon_{\mathrm{rr}}(\delta) = (5/8)^\kappa$.
The quantities of \cref{thm:cdhz} satisfy $q_s \le \kappa\ell + \ell + \lceil N_\ell/2 \rceil \le 2^{30}$ (positions read: one per query and level, one symbol per round polynomial, and the final table), $q_{\mathsf{V}} \le \nu + q_s(\lceil\log_2 l_{\max}\rceil + 1) \le 2^{36}$, and $q_{\widehat{\mathsf{CR}}} \le (M/2)\log_2 M \le 2^{36}$.
For an adversary making $Q = 2^{64}$ queries, \cref{thm:cdhz} then gives $\varepsilon_{\mathrm{ext}}(2^{64}) < 2^{-31.8}$: the first term is $2^{-31.84}$ and the other terms together are below $2^{-51}$, for every $n \le 30$ and $\ell \in \iv{1}{n}$ (computed exactly in rational arithmetic).
These are the figures of KBFold: the dominant term $320\, Q^2 \varepsilon_{\mathrm{rr}}(\delta)$ depends only on $(1-\delta)^\kappa$, and the smaller number of positions read does not change the other terms at this precision.
\end{remark}

\begin{remark}[Parameters]\label{rem:params}
For $\rho = 1/4$ and $\delta = (1-\rho)/2 = 3/8$, one query passes with probability at most $5/8$; $\kappa = 148$ queries give $(5/8)^{\kappa} < 2^{-100}$.
With $M = 2^{24}$ and $|\F| = p^2 > 2^{127.9}$ (the quadratic extension of \cref{sec:impl}), $\varepsilon_{\mathrm{fold}} + \ell/|\F| < 2^{-103}$, so the error of \cref{thm:sound} is below $2^{-100.2}$, also for $\Pieval^{\mathcal{J}}$ by \cref{thm:skip}(1).
These unique-decoding parameters are conservative; list-decoding analyses of FRI-type folding allow fewer queries (\cref{sec:concl}).
\end{remark}

\begin{remark}[Scope]\label{rem:fs}
\begin{itemize}
  \item The analysis concerns one commitment opened at one point $z$ per proof. Reusing a commitment across several openings, batching, or choosing $z$ after the commitment inside a larger Fiat--Shamir transcript is not analysed here.
  \item \Cref{cor:pq} applies to the compiler $\mathsf{BCS}$ of~\cite[Construction~11.7]{CDHZ25}, with salted Merkle trees and the challenge derivation of KBFold (\KB{\S3.5}).
  \item For $\ell = 0$ the same argument applies with $\Sigma = \F$, the implicit instance $y = w_0$ and the Hamming distance, giving the single error $(1-\delta)^{\kappa}$, as in \KB{Remark~7.30}.
  \item The bounds also hold against classical adversaries in the random oracle model, with the same (quantum) extractor.
\end{itemize}
\end{remark}

\subsection{Committed levels}\label{sec:sound:skip}

We transfer the results of this section to the protocol $\Pieval^{\mathcal{J}}$ of \cref{sec:scheme:skip}, for any set $\mathcal{J} \subseteq \iv{0}{\ell-1}$ of committed levels with $0 \in \mathcal{J}$, where $1 \le \ell \le n$.
The \emph{augmentation} $\mathrm{aug}$ maps a transcript of $\Pieval^{\mathcal{J}}$ to one of $\Pieval$ by inserting, in round $j+1$ for $j \in \iv{1}{\ell-1} \setminus \mathcal{J}$, the word $w_j = \cfold_{\theta_j}(w_{j-1})$ computed from the transcript, as in \cref{lem:skip}; it is computable in time polynomial in $M$, and the inserted word depends only on the part of the transcript before round $j+1$.
The IOR $\Pieval^{\mathcal{J},\Sigma}$, with the alphabet $\Sigma_{\mathcal{J}}$ of fibres and cosets, the permuted implicit instance $\varpi(\mathrm{fib}_0(w_0))$, the relation $\mathcal{R}^{\mathcal{J}}_{\mathrm{eval}}$ and the map $\mathrm{aug}^{\Sigma}$ on transcripts over $\Sigma_{\mathcal{J}}$, is defined as in \KB{\S7.8}, with $\cfold$ in place of the kernel fold and $P_f$ in place of $\tilde f$; the round polynomials are written over $\F^2$ as in \cref{sec:sound:pq}.

\begin{theorem}[Committed levels]\label{thm:skip}
Let $\mathcal{J} \subseteq \iv{0}{\ell-1}$ with $0 \in \mathcal{J}$.
\begin{enumerate}
  \item \Cref{thm:sound}, \cref{cor:binding} and \cref{cor:sum} hold for $\Pieval^{\mathcal{J}}$ with the same errors.
  \item $\Pieval^{\mathcal{J}}$ is round-by-round knowledge sound for $\mathcal{R}^{\delta}_{\mathrm{eval}}$ with the errors of \cref{thm:rbr}.
  \item $\Pieval^{\mathcal{J},\Sigma}$ is an IOR from $\mathcal{R}^{\mathcal{J}}_{\mathrm{eval}}$ to $\mathcal{R}_{\mathrm{acc}}$ with the relaxed round-by-round knowledge soundness errors of \cref{thm:rrbr}.
  \item \Cref{cor:pq} holds for $\mathsf{BCS}[\Pieval^{\mathcal{J},\Sigma}]$, with $y$ read through $\varpi^{-1}$, with the same $\nu$, $\beta_{\min}$ and $\varepsilon_{\mathrm{rr}}(\delta)$, with $l_{\max} = M/2$, and with $q_s \le \kappa\,(2^{d_0 - 1} + |\mathcal{J}| - 1) + \ell + \lceil N_\ell/2 \rceil$ positions read.
\end{enumerate}
\end{theorem}

\begin{proof}
(1) By \cref{lem:skip}, for every instance the acceptance probability of a prover $\tilde{\mathsf{P}}$ of $\Pieval^{\mathcal{J}}$ is at most that of the prover $\tilde{\mathsf{P}}'$ of $\Pieval$, over the same challenges. \Cref{thm:sound,cor:binding,cor:sum} quantify over all provers of $\Pieval$, and apply to $\tilde{\mathsf{P}}'$.

(2) Let $\Dd$ and $\mathsf{Ext}$ be as in \cref{thm:rbr}. Declare a partial transcript $\tau$ of $\Pieval^{\mathcal{J}}$ doomed if $\mathrm{aug}(\tau) \in \Dd$, and let $\mathsf{Ext}^{\mathcal{J}}(\tau, \mathsf{m}) = \mathsf{Ext}(\mathrm{aug}(\tau), \mathsf{m}')$, where $\mathsf{m}'$ is $\mathsf{m}$ together with the word that $\mathrm{aug}$ inserts in that round, if any.
Condition~1 of \cref{def:rbr-knowledge} holds because $\mathrm{aug}$ does not change the empty transcript.
For condition~2, $\mathrm{aug}(\tau, \mathsf{m}, \mathsf{c}) = (\mathrm{aug}(\tau), \mathsf{m}', \mathsf{c})$, because the inserted word depends only on $\mathrm{aug}(\tau)$; so the probability in condition~2 for $\Pieval^{\mathcal{J}}$ at $(\tau, \mathsf{m})$ equals the one for $\Pieval$ at $(\mathrm{aug}(\tau), \mathsf{m}')$, and the conclusion is the same.
For condition~3, if the verifier of $\Pieval^{\mathcal{J}}$ accepted a doomed full transcript $\tau$, the verifier of $\Pieval$ would accept the doomed transcript $\mathrm{aug}(\tau)$, by \cref{lem:skip}.

(3) and (4). The proofs of \KB{Theorem~7.31(3)--(4)} use only \KB{Lemma~6.5}, the form of $\mathrm{aug}^{\Sigma}$ and \KB{Theorem~7.25}; they apply verbatim with \cref{lem:skip} and \cref{thm:rrbr}.
The count of positions read differs from KBFold only in the round polynomials, which take one symbol per round instead of two.
\end{proof}

\begin{remark}[Parameters with committed levels]\label{rem:pq-skip}
For $\mathcal{J} = \{0, k, 2k, \dots\} \cap \iv{0}{\ell-1}$ with $k = 4$ and the parameters of \cref{rem:pq-params}, $q_s \le \kappa\,(2^{\min(4,\ell)-1} + \lceil \ell/4 \rceil - 1) + \ell + \lceil N_\ell/2 \rceil \le 2^{30}$, and the other quantities of \cref{rem:pq-params} are unchanged or smaller; so $\varepsilon_{\mathrm{ext}}(2^{64}) < 2^{-31.8}$ holds for $\mathsf{BCS}[\Pieval^{\mathcal{J},\Sigma}]$.
The implementation encodes the openings of one committed word by a merged authentication path (\cref{sec:impl:merkle}); as in \KB{Remark~7.32}, an adversary that outputs merged paths is turned into one that outputs separate paths with at most $q_s(\lceil\log_2 l_{\max}\rceil + 1) \le 2^{35}$ additional oracle queries, which does not change the bound at $Q = 2^{64}$ to the stated precision.
\end{remark}

\section{Implementation}\label{sec:impl}

We implemented $\Pieval^{\mathcal{J}}$ and $\Pisum$, compiled as in \cref{sec:proofs:bcs}, as a Rust library, \texttt{sumfri}, in the subdirectory \texttt{sumfri/rust} of the KBFold repository.
It depends on the \texttt{kbfold} crate (\KB{\S8}), from which it takes the field arithmetic, the NTT, the salted Merkle trees with merged openings, and the Fiat--Shamir transcript with the challenge maps $\phi$ and $\mathrm{pos}$; it adds the encoding, the restriction rounds, the classical fold, the prover and the verifier (about $700$ lines), the quotient baseline of \cref{sec:impl:quotient} (about $400$ lines), and about $500$ lines of tests and benchmarks.
Its only other dependencies are BLAKE3 and \texttt{getrandom}, as for \texttt{kbfold}.
This section records what differs from KBFold; \cref{sec:exp} reports measurements.

\paragraph{Index conventions.}
As in \KB{\S8}: bit $a_k$ of the paper is bit $k-1$ of the code, so that $a = a_1 + 2a'$ with $a_1$ the least significant bit; the coordinate $z_k$, the challenge $\theta_k$ and the round message of round $k$ are entries $k-1$ of their arrays.
A table $f$ with $n$ variables is an array of length $N$, and $f(a)$ is the coefficient of $X^a$ in $V_f$.

\subsection{Fields and domains}\label{sec:impl:fields}

Tables and $w_0$ live over the Goldilocks field $\mathbb{F}_p$, $p = 2^{64} - 2^{32} + 1$; the challenges $\theta_j$, the point $z$, the value $v$, the restricted tables $f_j$ ($j \ge 1$) and the folded words $w_j$ ($j \ge 1$) live in the quadratic extension $\K = \mathbb{F}_p[\iota]/(\iota^2 - 7)$ for the classical parameters, or in the quartic extension $\mathbb{F}_p[\iota]/(\iota^4 - 7)$ for the post-quantum parameters of \cref{rem:pq-params}.
The smooth domain is $L = \langle \omega \rangle$ with $\omega = 7^{(p-1)/M}$, stored in natural order; the fibre of index $i$ of $L_j$ is the pair of positions $(i, i + M_{j+1})$.
These are the choices of \KB{\S8.1--8.2}, unchanged.

\subsection{Encoding, rounds and folding}\label{sec:impl:enc}

\paragraph{Commit.}
$\Enc(f) = \ev_L(V_f)$ (\cref{def:encoding}) is computed by zero-padding the table to length $M$ and one radix-$2$ NTT. Since $V_f$ has coefficient vector $f$, there is no transform before the NTT; KBFold applies the kernel-to-monomial transform there ($\tfrac n2 N$ additions), and the coefficient-form baseline of KBFold applies M\"obius inversion (as many additions).

\paragraph{Round messages.}
The prover keeps the current table $f_{j-1}$ as an array over $\F$.
In round $j$ it sends only the linear coefficient $b_j$ of $s_j(T) = a_j + b_j T$ (\cref{rem:compress}); the verifier sets $a_j = v_{j-1} - z_j b_j$, so the restriction check holds by construction.
By \eqref{eq:honest-sj} and \cref{def:coeff-poly}, $b_j = \sum_{b'} f_{j-1}(2b'+1)\, \tau_j(b')$ with $\tau_j(b') = m_{b'}(z_{>j})$.
All the tensors $\tau_j$ are read from the single table $\tau_0 = (m_a(z))_{a}$ of length $N$: $\tau_j(b') = \tau_0(2^j b')$, since the low $j$ bits of $2^j b'$ vanish.
$\tau_0$ is built in $N$ multiplications, by doubling its length once per coordinate.
For $\Pisum$, $\tau_0 = \mathbf 1$; the code detects $z = \mathbf 1$ and replaces the inner products by sums.
After the challenge, the table is restricted in place, $f_j(b') = f_{j-1}(2b') + \theta_j f_{j-1}(2b'+1)$ (\cref{def:cres}): one multiplication per output entry, against two tables to restrict and three products per pair for the degree-$2$ sumcheck of KBFold (\KB{\S8.2}).

\paragraph{Fold.}
Prover and verifier compute $\cfold_\theta(w)(\xi^2) = a_e + \theta a_o$ with $a_e = (a + a')/2$ and $a_o = (a - a')/(2\xi)$, $a = w(\xi)$, $a' = w(-\xi)$ (\cref{lem:cword}(3)); the inverses $(2\xi)^{-1}$ come from one table of $\omega^{-i}$ shared by all levels, as in KBFold.
The same function serves the prover on whole words and the verifier on opened cosets.

\paragraph{Final message and closure.}
The verifier evaluates $G = V_g$ at the points $\xi_\ell$ by Horner's rule on $g$ directly: the final table is the coefficient vector of $G$, so the transform $C^{\otimes(n-\ell)}$ of \KB{\S8.3} disappears.
The closure value $P_g(z_{\ell+1},\dots,z_n)$ is computed by $n - \ell$ successive coefficient restrictions, in $O(N_\ell)$ operations; for $\Pisum$ it is $\sum_b g(b)$.

\subsection{Commitments and transcript}\label{sec:impl:merkle}

Committed levels, leaves, salts, openings, merged paths and local folds are those of \KB{\S8.3}, with $\mathcal{J} = \{0, k, 2k, \dots\}$, $k = 4$ by default.
The transcript has its own domain separator and absorbs the parameters $(n, R, \ell, \kappa)$, the salt length, the extension degree $e$ and $k$, the root of $w_0$, the point $z$ and the value $v$; in round $j$ it absorbs $b_j$, then the root of $w_{j-1}$ if $j-1 \in \mathcal{J}$, then a fresh round salt, and derives $\theta_j$ with the map $\phi$; in round $\ell+1$ it absorbs $g$ and a round salt, and derives the query points with $\mathrm{pos}$.
A proof consists of the $\ell$ values $b_j$, the roots, the $\ell+1$ round salts, $g$, and the openings of the committed words; it differs from a KBFold proof in the round messages, $\ell$ field elements instead of $3\ell$.

\subsection{The quotient baseline}\label{sec:impl:quotient}

For \cref{sec:exp}, the library also implements the standard univariate way to prove $\sum_a f(a) = v$ for the same commitment polynomial $V_f$: open $V_f$ at $X = 1$ with the quotient $q = (V_f - v)/(X-1)$, of degree less than $N - 1$, and test $q$ with FRI.
Since $1 \in L$, this baseline commits to $V_f$ on the coset $cL$, $c = 7$; the verifier computes the values of $q$ at the opened points from those of $V_f$ (with one batched inversion of the denominators $\xi - 1$), and the folding, committed levels, trees and transcript are those of $\Pisum$.
The final message is the coefficient vector of the folded quotient, $\mathrm{cres}_\theta(q)$ by \cref{thm:dictionary}, obtained from the coefficients $q_k = \sum_{a > k} f(a)$; there are no round messages.
This is the protocol that \cref{rem:deepfold} contrasts with $\Pisum$: the claim enters through a quotient word instead of through the shared challenges.

\subsection{Testing}\label{sec:impl:tests}

The test suite of \texttt{sumfri} checks:
\begin{itemize}
  \item \emph{Completeness.} Honest proofs of $\Pieval$ at random points and of $\Pisum$ verify for both extensions, $n \in \{1, 3, 6, 8, 10, 11\}$, several $\ell$ including $0$ and $n$, $k \in \{1, 2, 3, 4\}$ and $\rho \in \{1/2, 1/4, 1/8\}$, and the returned value equals $P_f(z)$, respectively $\sum_a f(a)$, computed independently from the definition.
  \item \emph{The restriction dictionary.} For $n \le 8$ and random $\theta \in \K^n$: $n$ restrictions give $P_f(\theta)$; $n$ classical folds of $\Enc(f)$ on a domain of size $4N$ give the constant word $P_f(\theta)$ (\cref{cor:cword-chain}); and $V_f(\zeta) = P_f(\zeta, \zeta^2, \dots)$ (\cref{lem:kronecker}).
  \item \emph{\Cref{ex:f17}}, step by step over $\mathbb{F}_{17}$.
  \item \emph{Tampering.} Proofs are rejected after changing the claimed value, one round message, one entry of $g$, one opened value, one round salt, the length of a list, or the commitment root.
  \item \emph{A consistent cheater.} A prover that commits honestly, claims $v + 1$, sends the honest $b_j$ (so that every restriction check holds), folds honestly and modifies $g(0)$ so that the closure check holds, is rejected in every tested configuration, at random points and at $z = \mathbf 1$, and always in the query phase, as case~2 of \cref{thm:sound} predicts.
  \item \emph{The quotient baseline}: completeness and rejection of a wrong value.
\end{itemize}

\section{Experiments}\label{sec:exp}

This section measures the implementation of \cref{sec:impl} against three schemes that prove the same statement with the same commitment machinery, and against an independent implementation.
For each measurement, every scheme runs on the same machine with the same field, NTT, Merkle trees, transcript, rate, number of queries and committed levels; they differ only in the protocol.

\subsection{Setup}\label{sec:exp:setup}

\paragraph{Machine and code.}
All runs use a cloud virtual machine with an Intel Xeon processor at $2.10$\,GHz ($2$ vCPUs, $7$\,GB RAM), the machine class of \KB{\S9}, Rust $1.97.0$, release profile with fat LTO and one codegen unit, and no \texttt{target-cpu} flag.
The libraries are \texttt{sumfri}~\sfversion{} and \texttt{kbfold}~0.3.0; the script \texttt{sumfri/rust/bench/run\_all.sh} runs every measurement of this section and writes the data of the tables to \texttt{sumfri/rust/results/}.
Runs are single-threaded.
Times are medians over $11$ runs for $n \le 16$, $5$ for $n \in \{18, 20\}$ and $3$ for $n = 22$; verification times are medians over $21$ runs.
The machine is shared; we repeated the measurements of \cref{tab:sum,tab:pq} in a second, independent run, whose figures are in the repository and lead to the same conclusions.

\paragraph{Parameters.}
The \emph{classical parameters} are those of \KB{\S9.1}: rate $\rho = 1/4$, $\ell = n - 4$ folding rounds (final table of $16$ elements), $\kappa = 148$ queries, the quadratic extension $\K$, $32$-byte salts, and committed levels $\mathcal{J} = \{0, 4, 8, \dots\}$.
By \cref{thm:sound,thm:skip} with $\delta = 3/8$, the soundness error is below $2^{-100}$ (\cref{rem:params}).
The \emph{post-quantum parameters} are those of \cref{rem:pq-params}: the same rate, $\ell$ and $\mathcal{J}$, $\kappa = 248$ queries and the quartic extension.
Tables are uniformly random over $\mathbb{F}_p$; evaluation points are uniformly random in $\F^n$.
Proof sizes are the lengths of the canonical encodings; they depend slightly on the query points, through the number of distinct opened leaves.

\paragraph{Schemes.}
\begin{itemize}
  \item \emph{Sum-FRI}: $\Pisum$, and $\Pieval$ at a random point (\cref{sec:impl}).
  \item \emph{KBFold}: the scheme of~\cite{Mkh26} (\texttt{kbfold}, kernel encoding), which proves the sum of its committed table as $2^n \tilde f(\tfrac12, \dots, \tfrac12)$ with a degree-$2$ sumcheck, and $\tilde f(z)$ at a random point.
  \item \emph{Coefficient form with sumcheck}: the coefficient-form mode of \texttt{kbfold} (\KB{\S9.1}), which commits to the Kronecker embedding of the coefficients of $\tilde f$ (after M\"obius inversion) and proves $\tilde f(z)$ with a degree-$2$ sumcheck and the classical fold, as in Reed--Solomon instantiations of BaseFold. For the sum, $z = (\tfrac12, \dots, \tfrac12)$.
  \item \emph{Quotient}: the opening of $V_f$ at $X = 1$ with a quotient and FRI (\cref{sec:impl:quotient}); it proves the sum only.
\end{itemize}
In all four schemes the commitment is one NTT of size $M$ and one salted tree of $M/2$ leaves, and the prover commits to the same number of folded words; the schemes differ in the transform before the NTT (none, kernel, M\"obius, none on a coset), in the round messages and the work that produces them, and, for the quotient, in the computation of the quotient word.
Every scheme commits to a table of $N$ field elements and proves the sum of its entries, so the four rows of \cref{tab:sum} prove the same statement about the same data.
For evaluations, Sum-FRI proves $P_f(z)$ and the other schemes $\tilde f(z)$, two different statements of the same cost class.

\subsection{The sum}\label{sec:exp:sum}

\begin{table}[t]
\centering
\small
\begin{tabular}{r r rrrr r r}
\toprule
 & Commit & \multicolumn{4}{c}{Open (ms)} & Verify & Proof \\
$n$ & (ms) & Sum-FRI & KBFold & coeff.+sc. & quotient & (ms) & (KiB) \\
\midrule
12 & 1.6 & 0.47 & 0.73 & 0.70 & 0.80 & 0.45 & 79 \\
14 & 7.2 & 2.1 & 3.2 & 3.0 & 3.1 & 0.60 & 115 \\
16 & 26.7 & 7.7 & 11.1 & 10.8 & 13.8 & 0.57 & 149 \\
18 & 136.4 & 31.9 & 43.6 & 45.3 & 61.5 & 0.90 & 193 \\
20 & 589.3 & 144.5 & 246.4 & 219.3 & 322.3 & 1.06 & 237 \\
22 & 3021.5 & 782.4 & 1150.1 & 1113.8 & 1474.4 & 1.32 & 288 \\
\bottomrule
\end{tabular}
\caption{Proving $\sum_a f(a)$ for a committed table of $N = 2^n$ elements, classical parameters, single-threaded. Commit, verify and proof size are those of Sum-FRI; the other three schemes are within the run-to-run variation of these figures, except the verifier of the quotient scheme, which is $6$--$56\%$ slower (\cref{sec:exp:sum}). The verification time at $n = 20$ is from the second run (the first gave $2.2$\,ms for Sum-FRI and $1.1$\,ms for the other schemes).}
\label{tab:sum}
\end{table}

\Cref{tab:sum} reports the four schemes on the sum.
The opening phase of Sum-FRI is $27$--$41\%$ faster than that of KBFold and $27$--$36\%$ faster than the coefficient-form scheme with sumcheck, over both runs and all sizes; it is $33$--$55\%$ faster than the quotient opening.
Commit times agree within the run-to-run variation (the four commitments perform the same NTT and the same tree; the transforms of KBFold and of the coefficient-form scheme cost $\tfrac n2 N$ additions, invisible next to the NTT), and so do proof sizes and the verification times of the two sumcheck schemes, which are dominated by the $\kappa$ authentication paths and opened cosets; the quotient verifier is $6$--$56\%$ slower, since it recomputes the quotient at the opened points.
The round messages, $\ell$ field elements for Sum-FRI and $3\ell$ for the sumcheck schemes, differ by at most $0.6$\,KiB.

The saving is the one predicted in \cref{sec:scheme:costs}.
With $z = \mathbf 1$, a restriction round of Sum-FRI costs one multiplication per entry of the table (the restriction $\mathrm{cres}_\theta$) and no multiplication for the round message (two sums), while a sumcheck round restricts two tables (the table and the equality table) and computes three products per pair.
The folds, the trees of the committed words and the openings are identical in the three schemes, and they account for the remaining time.
The quotient opening is slowest: it computes the quotient word on the whole domain, with one batched inversion of $M$ elements, before folding it.

\subsection{Evaluations}\label{sec:exp:eval}

\begin{table}[t]
\centering
\small
\begin{tabular}{r rrr r}
\toprule
 & \multicolumn{3}{c}{Open (ms)} & Proof \\
$n$ & Sum-FRI, $P_f(z)$ & KBFold, $\tilde f(z)$ & coeff.+sc., $\tilde f(z)$ & (KiB) \\
\midrule
12 & 0.75 & 0.76 & 0.75 & 79 \\
14 & 2.7 & 3.3 & 2.8 & 115 \\
16 & 12.1 & 14.1 & 13.6 & 148 \\
18 & 39.7 & 47.8 & 46.7 & 194 \\
20 & 219.3 & 256.8 & 264.0 & 236 \\
22 & 1064.7 & 1115.9 & 1123.8 & 290 \\
\bottomrule
\end{tabular}
\caption{Evaluation at a random point, classical parameters, single-threaded. Commit and verify times and proof sizes agree with \cref{tab:sum} within the run-to-run variation.}
\label{tab:eval}
\end{table}

At a random point (\cref{tab:eval}), the opening of Sum-FRI is $0$--$17\%$ faster than both sumcheck schemes.
The restriction round now needs the tensor $\tau_j$ (\cref{sec:impl:enc}): two inner products with $\tau_j$ per round, plus $N$ multiplications once to build $\tau_0$, against three products per pair and the restriction of the equality table for the sumcheck.
The gap is smaller than for the sum, and at $n = 22$ it is within the run-to-run variation.

\subsection{Post-quantum parameters}\label{sec:exp:pq}

\begin{table}[t]
\centering
\small
\begin{tabular}{r rr rr rr}
\toprule
 & \multicolumn{2}{c}{Open, sum (ms)} & \multicolumn{2}{c}{Open, evaluation (ms)} & Verify & Proof \\
$n$ & Sum-FRI & KBFold & Sum-FRI & KBFold & (ms) & (KiB) \\
\midrule
12 & 0.87 & 1.42 & 1.03 & 1.42 & 1.07 & 131 \\
14 & 3.3 & 7.1 & 5.1 & 7.6 & 1.38 & 215 \\
16 & 14.1 & 29.0 & 21.6 & 26.5 & 1.68 & 295 \\
18 & 58.2 & 116.1 & 88.1 & 121.9 & 2.00 & 390 \\
20 & 269.3 & 626.5 & 468.9 & 605.9 & 3.30 & 498 \\
22 & 1443.5 & 2469.5 & 1971.6 & 2536.5 & 2.87 & 619 \\
\bottomrule
\end{tabular}
\caption{Post-quantum parameters (quartic extension, $\kappa = 248$), single-threaded. Verify and proof size are those of Sum-FRI for the sum. Commit times are those of \cref{tab:sum} within the run-to-run variation: the commitment does not depend on the extension.}
\label{tab:pq}
\end{table}

With the post-quantum parameters (\cref{tab:pq}), every operation of the rounds is in the quartic extension, built over $\K$ by Karatsuba (three multiplications in $\K$ per product).
Sum-FRI opens the sum $34$--$57\%$ faster than KBFold over both runs, and evaluations $19$--$33\%$ faster.
Proofs are $1.7$--$2.2$ times larger than with the classical parameters, as in \KB{\S9.3}: $100$ more queries, each opening elements of $\mathbb{F}_{p^4}$.

\subsection{Where the time goes}\label{sec:exp:breakdown}

At $n = 20$ with the classical parameters, the commitment spends $251$\,ms in the NTT and $490$\,ms in the salted tree of $M/2 = 2^{21}$ leaves (medians of separate runs of each step, whose sum exceeds the median commit time of the same run, $657$\,ms, by the run-to-run variation).
The opening of the sum takes $161$\,ms and that of an evaluation $203$\,ms: the $42$\,ms difference is the cost of the tensor $\tau_0$ and of the products with it.
For all four schemes the commitment dominates the prover, and its cost is shared: Sum-FRI changes only the opening.

\subsection{An independent reference point}\label{sec:exp:whir}

\begin{table}[t]
\centering
\small
\begin{tabular}{r rrr rrr rrr}
\toprule
 & \multicolumn{3}{c}{Sum-FRI, classical} & \multicolumn{3}{c}{WHIR, unique, $k = 1$} & \multicolumn{3}{c}{WHIR, Johnson, $k = 4$} \\
$n$ & prover & verify & proof & prover & verify & proof & prover & verify & proof \\
\midrule
12 & 2.1 & 0.45 & 79 & 18.4 & 1.40 & 131 & 4.6 & 0.38 & 45 \\
14 & 9.2 & 0.60 & 115 & 87.7 & 2.30 & 199 & 19.2 & 0.53 & 56 \\
16 & 34.4 & 0.57 & 149 & 335.1 & 2.80 & 285 & 71.2 & 0.66 & 75 \\
18 & 168.3 & 0.90 & 193 & 1318 & 3.90 & 392 & 340.6 & 0.48 & 88 \\
20 & 733.8 & 1.06 & 237 & 6183 & 3.40 & 513 & 2470 & 0.54 & 108 \\
22 & 3803.9 & 1.32 & 288 & 27000 & 5.50 & 657 & 9700 & 0.54 & 121 \\
\bottomrule
\end{tabular}
\caption{Sum-FRI ($\Pisum$, prover $=$ commit $+$ open) and the WHIR implementation, same machine, single-threaded, Goldilocks with its quadratic extension, rate $1/4$, $100$ bits of security, no proof of work. Times in ms, proof sizes in KiB. WHIR prints times of one second or more with a resolution of $0.1$\,s; its figures are those of the median, by commit time, of three invocations, each with \texttt{--reps 21}. The Sum-FRI verification time at $n = 20$ is from the second run, as in \cref{tab:sum}.}
\label{tab:whir}
\end{table}

As in \KB{\S9.8}, we run the WHIR implementation~\cite{ACFY24}\footnote{Repository \url{https://github.com/WizardOfMenlo/whir}, commit \texttt{c03a4a5}, the version used in KBFold.} on the same machine with the same field, rate and security level, single-threaded and without proof of work, in two configurations: one variable per round with the unique-decoding analysis, and four variables per round with its analysis up to the Johnson bound (\cref{tab:whir}).
WHIR proves a multilinear evaluation, and the two code bases differ in many engineering choices, so \cref{tab:whir} is a reference point, not a comparison of protocols.
The Sum-FRI prover is $7$--$10$ times faster than WHIR with one variable per round, with proofs $1.7$--$2.3$ times smaller and a verifier $3$--$5$ times faster.
Against WHIR's default configuration, the Sum-FRI prover is $2$--$3.4$ times faster, while WHIR's proofs are $1.8$--$2.4$ times smaller; the verification times are comparable for $n \le 16$, and WHIR's verifier is $1.9$--$2.4$ times faster for $n \ge 18$. As in KBFold, the difference in proof size and verification time comes from the number of queries, which WHIR reduces by its analysis up to the Johnson bound.

\section{Conclusion and Future Work}\label{sec:concl}

Classical FRI folding of the Kronecker encoding $V_f$ restricts the coefficient polynomial $P_f$ in its first variable (\cref{thm:dictionary}).
When the folding challenges also drive the claim, an evaluation $P_f(z)$, and in particular the hypercube sum $\sum_a f(a) = P_f(\mathbf 1)$, is reduced to a claim on the final table by one affine round polynomial per round, with no equality table and no change of basis.
This round structure is DeepFold's~\cite{GLHQTZ25}.
We gave a complete proof of its soundness, evaluation binding, round-by-round knowledge soundness and, through the theorem of Chiesa, Di, Hu and Zheng~\cite{CDHZ25}, post-quantum knowledge soundness of the compiled scheme for a single opening, in the unique-decoding regime.
The proof follows KBFold: correlated agreement enters through the line $w_e + \theta w_o$, and the sumcheck step becomes a restriction step that costs one bad challenge per round.
All of it except the quantum compiler, the Merkle-tree statements, the decoding algorithm and running times is checked in Lean, against the single axiom of correlated agreement.
The implementation shares everything with KBFold except the protocol, and opening a sum is $27$--$41\%$ faster than with KBFold or with a coefficient-form sumcheck at equal commitment cost, verification time and proof size, and, with post-quantum parameters, $34$--$57\%$ faster than with KBFold.

\subsection{Future work}\label{sec:concl:future}

\paragraph{Soundness beyond unique decoding.}
\begin{itemize}
  \item \emph{List decoding.} Extend \cref{thm:sound,thm:rbr} to proximity parameters up to the Johnson bound. The classical word fold is directly a Reed--Solomon line (\cref{lem:cword}(2)), so correlated agreement beyond unique decoding~\cite{BCIKS23}, or its subcode form~\cite{Hab24}, applies to it; the main work is to replace the unique codewords of the codeword chain by lists, and to carry the restriction step through the list. This would reduce the number of queries, which dominates the proof size and the verification time in \cref{tab:whir}.
  \item \emph{Out-of-domain samples.} DeepFold~\cite{GLHQTZ25} adds out-of-domain evaluations to reach list decoding; with them, the analysis of \cref{sec:sound} would have to account for the additional claims on the folded polynomials.
\end{itemize}

\paragraph{Batching and composition.}
\begin{itemize}
  \item \emph{Several sums.} Sums of several committed tables, or of one table against several points, by a random linear combination before the first round; soundness by correlated agreement applied to the combination, as for batched FRI.
  \item \emph{Linear claims.} $\Pieval$ proves the inner product of the committed table with the tensor $(m_a(z))_a$. Other structured public vectors, and the use of $\Pisum$ as the final step of protocols that reduce their claims to a sum of one committed table, are natural next steps.
\end{itemize}

\paragraph{Zero knowledge.}
Masking the committed polynomial with $X^N R$ for a random $R$ of small degree, masking the round polynomials, and masking the folded words with a random codeword, as proposed in \KB{\S10.1}; the restriction dictionary has to be extended to the masked polynomial, and the view shown to be simulable.

\paragraph{Formal verification.}
Replace the axiom \texttt{bciks\_unique} by a proof of correlated agreement in the unique-decoding regime; and link the Rust implementation to the Lean specification.

\subsection{Artefacts}\label{sec:concl:artefacts}
The paper, the Rust library \texttt{sumfri}, the Lean library \texttt{SumFRI} and the benchmark data are in the KBFold repository, \url{https://github.com/qoosmo/kbfold}, under the MIT or Apache-2.0 licences (code) and CC~BY~4.0 (paper).
\begin{itemize}
  \item \emph{Tests} (\cref{sec:impl:tests}): \texttt{cargo test --release} in \texttt{sumfri/rust/}.
  \item \emph{Tables of \cref{sec:exp}}: in \texttt{sumfri/rust/},\\ \texttt{bench/run\_all.sh <tag> 22 <whir\_dir>}, or \texttt{cargo run --release --example bench -- sum} (and \texttt{eval}, \texttt{pq}, \texttt{breakdown}); the recorded data are in \texttt{sumfri/rust/results/}.
  \item \emph{Lean} (Lean~4.23.0, Mathlib \texttt{v4.23.0}): \texttt{lake build SumFRI} in \texttt{lean/};\\ \texttt{lake env lean SumFRI/Audit.lean} prints the axioms of the main theorems.
\end{itemize}


\section*{Acknowledgments}
The author thanks the developers of Lean, Mathlib and the Rust ecosystem, on which the companion artefacts are built.
Generative AI tools (Anthropic's Claude) were used to assist with writing and editing the text, with the Rust code and the Lean proofs, and with checking computations.
The author checked all results and takes full responsibility for the content of the paper and the artefacts; every Lean proof is checked by the Lean kernel.

\bibliographystyle{alpha}
\bibliography{refs}

\appendix
\section{Notation}\label{app:notation}

\begin{center}
\small
\begin{longtable}{@{}p{0.27\textwidth} p{0.49\textwidth} >{\raggedright\arraybackslash}p{0.17\textwidth}@{}}
\toprule
Symbol & Meaning & Where \\
\midrule
\endhead
\multicolumn{3}{l}{\emph{Integers, bits and fields}} \\
$\iv{a}{c}$; $n$, $N = 2^n$ & integers $a \le i \le c$; number of variables; size of a table & \S\ref{sec:prelim:bits} \\
$a_k$, $a = a_1 + 2a'$ & bit $k$ of an index (least significant first); splitting off the first bit & \S\ref{sec:prelim:bits} \\
$\Fq$, $\F$, $\K$ & base field; challenge field $\F \supseteq \Fq$; extension field & \S\ref{sec:prelim:fields}, \S\ref{sec:scheme:params} \\
$\F[X]_{<d}$, $\coef{a}{P}$, $P_e$, $P_o$ & polynomials of degree $< d$; coefficient of $X^a$; even and odd parts & \Cref{lem:splits} \\
\midrule
\multicolumn{3}{l}{\emph{Tables and polynomials}} \\
$f$, $f(a)$ & table, $a \in \iv{0}{N-1}$ & \S\ref{sec:prelim:ml} \\
$m_a$, $\eq$, $\tilde f$ & canonical monomial; equality polynomial; multilinear extension & \Cref{def:bases,prop:mle} \\
$P_f$ & coefficient polynomial $\sum_a f(a)\, m_a$ & \Cref{def:coeff-poly} \\
$V_f$ & Kronecker encoding $\sum_a f(a) X^a$ & \Cref{def:vf} \\
$\mathbf 1$ & the point $(1, \dots, 1)$; $P_f(\mathbf 1) = V_f(1) = \sum_a f(a)$ & \eqref{eq:sum-is-eval}, \Cref{lem:kronecker} \\
$g_e$, $g_o$ & $g_e(b') = g(2b')$, $g_o(b') = g(2b'+1)$ & \Cref{def:round-poly} \\
$\mathrm{cres}_\theta(f)$, $\mathrm{cres}_{\theta_{\le j}}(f)$ & coefficient restriction $f_e + \theta f_o$; iterated & \Cref{def:cres} \\
$\res_r(f)$ & restriction of the multilinear extension (KBFold) & \Cref{def:restriction} \\
$U_f$ & kernel encoding of KBFold, $\sum_b f(b) K_b$ & \Cref{rem:kernel-relation} \\
\midrule
\multicolumn{3}{l}{\emph{Domains, codes and folds}} \\
$L$, $M$; $L_j$, $M_j$ & smooth domain and its order; $L^{2^j}$ and its order $M/2^j$ & \Cref{def:domain}, \S\ref{sec:scheme:params} \\
fibre $\{\xi, -\xi\}$ & the two square roots of $\xi^2$ & \Cref{def:fibre} \\
$\ev_L$, $\RS_\F[L,d]$, $\rho$, $\Delta$ & evaluation map; Reed--Solomon code; rate; relative Hamming distance & \Cref{def:rs} \\
$w_e$, $w_o$ & even and odd parts of a word & \eqref{eq:even-odd} \\
$\cfold_\theta(U)$, $\cfold_\theta(w)$ & classical fold $U_e + \theta U_o$ of a polynomial, $w_e + \theta w_o$ of a word & \Cref{def:cfold,lem:cword} \\
$\mathrm{pfold}_r$ & kernel fold of KBFold & \Cref{rem:kernel-relation} \\
$h_{g,y}$ & round polynomial $P_{g_e}(y) + T P_{g_o}(y)$ & \Cref{def:round-poly} \\
\midrule
\multicolumn{3}{l}{\emph{The scheme}} \\
$R$, $\ell$, $\kappa$ & $\rho = 2^{-R}$; number of folding rounds; number of queries & \S\ref{sec:scheme:params} \\
$N_j$, $\Cc_j$ & $N/2^j$; $\RS_\F[L_j, N_j]$ & \S\ref{sec:scheme:params} \\
$\Enc(f)$ & $\ev_L(V_f)$ & \Cref{def:encoding} \\
$\Pieval$, $\Pisum$ & evaluation protocol; sum protocol ($z = \mathbf 1$) & \S\ref{sec:scheme:protocol} \\
$z$, $v$, $z_{>j}$ & point and claimed value; $(z_{j+1}, \dots, z_n)$ & \S\ref{sec:scheme:protocol} \\
$s_j = a_j + b_j T$, $\theta_j$, $v_j$ & round polynomial; challenge; running claim $s_j(\theta_j)$ & \S\ref{sec:scheme:protocol} \\
$f_j$, $g$, $G$, $w_\ell$ & honest tables $\mathrm{cres}_{\theta_{\le j}}(f)$; final table; $V_g$; $\ev_{L_\ell}(G)$ & \S\ref{sec:scheme:protocol} \\
$\tau_j$ & monomial tensor $(m_{b'}(z_{>j}))_{b'}$ & \S\ref{sec:scheme:costs} \\
$\mathcal{J}$, $d_j$, $\Pieval^{\mathcal{J}}$ & committed levels; folds between committed levels; protocol with committed levels & \S\ref{sec:scheme:skip} \\
\midrule
\multicolumn{3}{l}{\emph{Soundness}} \\
$\delta$, $\delta^*$ & proximity parameter; $(1-\rho)/2$ & \S\ref{sec:sound}, \S\ref{sec:sound:pq} \\
$\Delta^{\mathrm{fib}}_j$ & fibre distance & \Cref{def:fibre-dist} \\
$\mathrm{dec}_j(w)$, $f^*$ & decoded codeword; decoded table & \Cref{def:decoded} \\
$\mathsf{Good}_j(w, \theta)$ & good challenge & \Cref{def:good} \\
$\Pp_j$, $\pi$ & passing sets; fraction of passing query points & \S\ref{sec:sound:thm} \\
$\varepsilon_{\mathrm{fold}}$ & $\sum_{j=1}^\ell M_j/|\F|$ & \Cref{thm:sound} \\
$\lambda[c]$, $\sigma_j(c)$ & coefficient vector of $P_c$; $P_{\lambda[c]}(z_{>j})$ & \S\ref{sec:sound:rbr} \\
$\hat w_j$, $\mathcal{B}_j$, $\mathcal{W}_j(c)$, $\mathrm{dens}_j$ & folded word; disagreement set; witness set; its density & \S\ref{sec:sound:rbr} \\
$\Dd$, $\mathsf{Ext}$ & doomed set; extractor & \Cref{def:state} \\
$\Sigma$, $\mathrm{fib}_j$, $\phi$, $\mathrm{pos}$, $\beta$ & alphabet $\F^2$; fibre string; challenge maps; challenge length & \S\ref{sec:sound:pq} \\
\bottomrule
\end{longtable}
\end{center}

\section{Lean formalisation}\label{app:lean}

The Lean~4 library \texttt{SumFRI} (Lean~4.23.0, Mathlib at the tag \texttt{v4.23.0}; about $3900$ lines in $11$ modules and an audit file, in the directory \texttt{lean/SumFRI} of the KBFold repository) is a second library of the KBFold Lean project.
It imports the KBFold library~\cite{Mkh26} and reuses its multilinear tables, smooth domains, Reed--Solomon codes, fibre distance, probability lemmas, generic round-by-round framework, fibre strings and sampling lemmas.
It contains no \texttt{sorry}. Its only axiom is KBFold's \texttt{bciks\_unique}, the correlated-agreement theorem (\cref{thm:bciks}); \texttt{\#print axioms} (file \texttt{SumFRI/Audit.lean}) shows that the restriction dictionary, completeness, the case $\ell = 0$ and the deterministic core of \cref{lem:cr-unique} do not use it, and that soundness, binding, round-by-round knowledge soundness and relaxed round-by-round knowledge soundness use it and the three standard axioms of Lean (\texttt{propext}, \texttt{Classical.choice}, \texttt{Quot.sound}) only.

\paragraph{Modelling.}
The modelling is that of KBFold (\KB{Appendix~B}).
Bit vectors are functions $\mathtt{Fin}\ n \to \mathtt{Bool}$, with bit $0$ of Lean the bit $a_1$ of the paper; one field $\F$ contains the smooth domain $L \le \F^\times$; probabilities are rationals computed by counting; rounds and challenges are indexed from $0$.
A deterministic prover is a family of messages indexed by the challenges, with a causality hypothesis and the hypothesis that round polynomials have degree at most $1$ (they are transmitted as two coefficients, or as one with \cref{rem:compress}).
The point $z$ is given by its first $\ell$ coordinates and its last $n - \ell$ coordinates, so that the table has $n = m + \ell$ variables with $m = n - \ell$.

\paragraph{Correspondence.}
\begin{center}
\small
\begin{longtable}{@{}>{\raggedright\arraybackslash}p{0.34\linewidth}>{\raggedright\arraybackslash}p{0.18\linewidth}>{\raggedright\arraybackslash}p{0.42\linewidth}@{}}
\toprule
Paper & Lean module & Main declarations \\
\midrule
\endhead
\Cref{def:coeff-poly,def:vf,lem:kronecker,def:cres,lem:cres,thm:dictionary}, \eqref{eq:sum-is-eval} & \texttt{Monomial} & \texttt{cpoly}, \texttt{vpoly}, \texttt{vpoly\_eval}, \texttt{vpoly\_eval\_one}, \texttt{cres}, \texttt{cpoly\_cres}, \texttt{cpoly\_cresSeq}, \texttt{cfold\_vpoly}, \texttt{cfoldSeq\_vpoly\_full}, \texttt{cpoly\_one} \\
\Cref{lem:cword}(4), \cref{def:round-poly,lem:round-poly} & \texttt{Monomial} & \texttt{cwfold\_ev}, \texttt{rnd0}, \texttt{rnd1}, \texttt{hround\_eq}, \texttt{hround\_one}, \texttt{affine\_agree} \\
\Cref{lem:cword}(3), \cref{lem:local-fold} & \texttt{LocalFold} & \texttt{cwfold\_sqPt}, \texttt{cfoldUp\_local} \\
\Cref{def:encoding,cor:cword-chain,lem:honest,thm:complete} & \texttt{Scheme} & \texttt{cEnc}, \texttt{cfoldUp\_Enc}, \texttt{hS}, \texttt{hS\_eval}, \texttt{hS\_zp}, \texttt{SumProver}, \texttt{completeness}, \texttt{completeness\_sum} \\
\Cref{lem:encoding,def:decoded}, \cref{lem:cword}(1), \cref{lem:far,lem:survive,def:good,lem:good,lem:step} & \texttt{Fibre} & \texttt{vpoly\_ccoeffs}, \texttt{cEnc\_injective}, \texttt{cdecTable}, \texttt{cwfold\_congr}, \texttt{cfar\_fold}, \texttt{csurvive}, \texttt{CGood}, \texttt{ncard\_not\_cgood\_le}, \texttt{cone\_step} \\
\Cref{lem:passing,lem:chain,thm:sound,cor:binding,cor:sum,prop:l0} & \texttt{Soundness} & \texttt{cPass\_card}, \texttt{cchain}, \texttt{cchain\_final}, \texttt{csoundness\_far}, \texttt{csoundness\_close}, \texttt{ceval\_binding}, \texttt{soundness\_sum}, \texttt{sum\_binding}, \texttt{cno\_folding} \\
\Cref{def:state,lem:rbr-step,lem:rbr-restr,thm:rbr,prop:rbr-l0} & \texttt{RBR} & \texttt{cNotDoomed}, \texttt{cExt}, \texttt{crbr\_step}, \texttt{ccoeffs\_cwfold}, \texttt{csRw\_zp}, \texttt{csRw\_eval}, \texttt{crbr\_round\_one}, \texttt{crbr\_round}, \texttt{crbr\_final}, \texttt{crbr\_l0} \\
\Cref{thm:rbr} (``Consequently''), \cref{rem:bgtz} & \texttt{RBRGeneric} & \texttt{crbr\_knowledge}, \texttt{crbr\_binding} \\
\Cref{def:kstate,lem:rbr-erasures,thm:rrbr}, \cref{lem:cr-unique} (core) & \texttt{NonInteractive} & \texttt{cKState}, \texttt{crbr\_stepE}, \texttt{crbr\_erasures\_one}, \texttt{crbr\_erasures\_round}, \texttt{crbr\_erasures\_final}, \texttt{crrbr\_round}, \texttt{crrbr\_final}, \texttt{ccr\_unique\_core} \\
\Cref{lem:skip}, \cref{thm:complete} for $\Pieval^{\mathcal{J}}$, \cref{thm:skip}(1) & \texttt{Skip} & \texttt{SumProver.augW}, \texttt{caccepts\_aug}, \texttt{completeness\_J}, \texttt{completeness\_sum\_J}, \texttt{csoundness\_far\_J}, \texttt{csoundness\_close\_J}, \texttt{ceval\_binding\_J}, \texttt{ceval\_binding'\_J}, \texttt{soundness\_sum\_J}, \texttt{sum\_binding\_J} \\
\Cref{thm:skip}(2) & \texttt{SkipRBR} & \texttt{augPT}, \texttt{augWord\_updR}, \texttt{crbr\_round\_one\_J}, \texttt{crbr\_round\_J}, \texttt{crbr\_final\_J} \\
\Cref{thm:skip}(3) & \texttt{Skip\-Non\-Interactive} & \texttt{augE}, \texttt{cNotDoomedE\_congr}, \texttt{cAcceptsSigJ\_aug}, \texttt{crrbr\_round\_J}, \texttt{crrbr\_final\_J} \\
\bottomrule
\end{longtable}
\end{center}

\paragraph{Reused from KBFold.}
\Cref{lem:sequential,lem:queries,lem:randomised,def:rbr-knowledge,lem:rbr-overall,lem:rbr-binding} (module \texttt{RoundByRound} and \texttt{Probability}), \cref{lem:power,def:fibre,def:rs,lem:metric,lem:rs-params,lem:split-words} (\texttt{Domain}), \cref{def:fibre-dist,lem:fibre-dist,lem:unique} (\texttt{Fibre}), \cref{lem:sampling} and the fibre strings of \cref{sec:sound:pq} (\texttt{NonInteractive}), \cref{def:cfold} and \cref{lem:cword}(2) (\texttt{WordFold}).

\paragraph{Not formalised.}\leavevmode\\
\Cref{thm:cdhz} and \cref{cor:pq}, which rest on the quantum analysis of~\cite{CDHZ25}; the Merkle-tree statements and the collision step of \cref{lem:cr-unique}; \cref{fact:decode} and \cref{lem:descent}; running times and operation counts (in particular the count of \cref{lem:local-fold}); item 4 of \cref{thm:skip}; \cref{sec:impl,sec:exp}.
The Lean verifier of $\Pieval^{\mathcal{J}}$ states each check with the classical fold of the augmented word, which the verifier computes from the opened coset by \cref{lem:local-fold} (formalised as \texttt{cfoldUp\_local}); items 2 and 3 of \cref{thm:skip} are stated, like items 1--3 of \cref{thm:rbr}, \cref{thm:rrbr} and, for condition~2 only, \cref{prop:rbr-l0}, for the per-round bounds, with the doomed set and the knowledge state function pulled back through the augmentation, and a committed coset word is represented by its fibre string; and \cref{lem:rbr-erasures,thm:rrbr} are stated for a model of transcripts by stages.

\paragraph{Findings of the formalisation.}
The formalisation found no false statement. As in KBFold, case~2 of \cref{thm:sound} holds without its hypothesis $\Delta^{\mathrm{fib}}_0(w_0, \Cc_0) \le \delta$, and the Lean statement is in this form.


\end{document}
